
\documentclass[11pt]{article}

\usepackage[margin=0.85in]{geometry}
\usepackage{amsmath,amssymb,mathtools}
\usepackage{enumitem}
\usepackage{xcolor}
\usepackage[most]{tcolorbox}
\usepackage{fancyhdr}
\usepackage{lastpage}
\usepackage{hyperref}
\usepackage{microtype}
\usepackage{lmodern}

%%%%%%%%%%%%%%%%%%%%%%%%%%%%%%%%%%%%%
% Colors
%%%%%%%%%%%%%%%%%%%%%%%%%%%%%%%%%%%%%

\definecolor{uwpurple}{HTML}{4B2E83}
\definecolor{uwgold}{HTML}{B7A57A}
\definecolor{softpurple}{HTML}{F4F1FA}

%%%%%%%%%%%%%%%%%%%%%%%%%%%%%%%%%%%%%
% Hyperlinks
%%%%%%%%%%%%%%%%%%%%%%%%%%%%%%%%%%%%%

\hypersetup{
    colorlinks=true,
    linkcolor=uwpurple,
    urlcolor=uwpurple,
    citecolor=uwpurple
}

%%%%%%%%%%%%%%%%%%%%%%%%%%%%%%%%%%%%%
% Page Style
%%%%%%%%%%%%%%%%%%%%%%%%%%%%%%%%%%%%%

\pagestyle{fancy}
\fancyhf{}
\lhead{\textbf{MATH 394: Probability I}}
\rhead{\textbf{Practice Problems}}
\cfoot{\thepage\ of \pageref{LastPage}}

\setlength{\headheight}{14pt}
\setlength{\parindent}{0pt}
\setlength{\parskip}{0.45em}

\setlist[enumerate]{
    itemsep=0.35em,
    topsep=0.35em
}

%%%%%%%%%%%%%%%%%%%%%%%%%%%%%%%%%%%%%
% Course Information
%%%%%%%%%%%%%%%%%%%%%%%%%%%%%%%%%%%%%

\newcommand{\course}{MATH 394: Probability I}
\newcommand{\documenttitle}{Question Bank with Solutions \\ (Midterm)}
\newcommand{\instructor}{Arman Jahangiri}
\newcommand{\department}{Department of Mathematics}
\newcommand{\university}{University of Washington}
\newcommand{\term}{Summer 2026}

%%%%%%%%%%%%%%%%%%%%%%%%%%%%%%%%%%%%%
% Problem Formatting
%%%%%%%%%%%%%%%%%%%%%%%%%%%%%%%%%%%%%

\newcounter{problem}

\newtcolorbox{problemheader}{
    colback=softpurple,
    colframe=uwpurple,
    boxrule=0.8pt,
    arc=2mm,
    left=2mm,
    right=2mm,
    top=1.5mm,
    bottom=1.5mm,
    before skip=1em,
    after skip=0.5em
}

\newenvironment{problem}[1][]{
    \refstepcounter{problem}
    \begin{problemheader}
        \textbf{Problem \theproblem%
        \if\relax\detokenize{#1}\relax
        \else. #1%
        \fi}
    \end{problemheader}
}{
    \vspace{0.5em}
}

%%%%%%%%%%%%%%%%%%%%%%%%%%%%%%%%%%%%%
% Solution Formatting
%%%%%%%%%%%%%%%%%%%%%%%%%%%%%%%%%%%%%

\newenvironment{solution}{
    \par
    \begin{tcolorbox}[
        colback=white,
        colframe=uwgold,
        boxrule=0.7pt,
        arc=2mm,
        left=3mm,
        right=3mm,
        top=2mm,
        bottom=2mm,
        breakable,
        before skip=0.5em,
        after skip=1em,
        title=\textbf{Solution},
        coltitle=uwpurple,
        fonttitle=\bfseries
    ]
}{
    \end{tcolorbox}
}

%%%%%%%%%%%%%%%%%%%%%%%%%%%%%%%%%%%%%
% Optional Source Command
%%%%%%%%%%%%%%%%%%%%%%%%%%%%%%%%%%%%%

\newcommand{\source}[1]{%
    \vspace{0.1cm}
    
    {\small\textit{Source: #1}}
}

%%%%%%%%%%%%%%%%%%%%%%%%%%%%%%%%%%%%%
% Optional Answer Boxes
%%%%%%%%%%%%%%%%%%%%%%%%%%%%%%%%%%%%%

\newcommand{\answerbox}{%
    \begin{flushright}
        \begin{tcolorbox}[
            width=0.36\textwidth,
            height=2cm,
            colback=white,
            colframe=uwpurple,
            boxrule=0.7pt,
            arc=1mm
        ]
            {\small\textbf{Final Answer}}
        \end{tcolorbox}
    \end{flushright}
}

\newcommand{\partanswerbox}[1]{%
    \begin{flushright}
        \begin{tcolorbox}[
            width=0.36\textwidth,
            height=2cm,
            colback=white,
            colframe=uwpurple,
            boxrule=0.7pt,
            arc=1mm
        ]
            {\small\textbf{Final Answer for Part #1}}
        \end{tcolorbox}
    \end{flushright}
}

%%%%%%%%%%%%%%%%%%%%%%%%%%%%%%%%%%%%%
% Mathematical Commands
%%%%%%%%%%%%%%%%%%%%%%%%%%%%%%%%%%%%%

\newcommand{\E}{\mathbb{E}}
\newcommand{\Pp}{\mathbb{P}}
\newcommand{\cC}{\mathcal{C}}

%%%%%%%%%%%%%%%%%%%%%%%%%%%%%%%%%%%%%
% Document
%%%%%%%%%%%%%%%%%%%%%%%%%%%%%%%%%%%%%

\begin{document}

%%%%%%%%%%%%%%%%%%%%%%%%%%%%%%%%%%%%%
% Title Page
%%%%%%%%%%%%%%%%%%%%%%%%%%%%%%%%%%%%%

\begin{titlepage}
    \thispagestyle{empty}

    \begin{center}

        \vspace*{1.2in}

        {\color{uwpurple}\rule{\textwidth}{1.5pt}}

        \vspace{0.45in}

        {\LARGE\bfseries\color{uwpurple}
        \documenttitle}

        \vspace{0.25in}

        {\Large\bfseries
        \course}

        \vspace{0.45in}

        {\color{uwgold}\rule{0.65\textwidth}{1.2pt}}

        \vfill

        \begin{tcolorbox}[
            width=0.72\textwidth,
            colback=softpurple,
            colframe=uwpurple,
            boxrule=0.8pt,
            arc=2mm,
            left=5mm,
            right=5mm,
            top=4mm,
            bottom=4mm
        ]
            \centering
            \large
            \textbf{Instructor}\\
            \instructor

            \vspace{0.3cm}

            \textbf{Term}\\
            \term
        \end{tcolorbox}

        \vfill

        {\large\bfseries\university}

        \vspace{0.1in}

        {\large\department}

        \vspace{0.6in}

        {\color{uwpurple}\rule{\textwidth}{1.5pt}}

    \end{center}
\end{titlepage}

\setcounter{page}{1}






\setcounter{problem}{0}



\begin{problem}
How many different seven-place license plates are possible if the first two places are for letters and the other five are for digits?
\begin{enumerate}[label=(\alph*)]
    \item Repetition is allowed.
    \item No letter or digit can be repeated on a single license plate.
\end{enumerate}
\end{problem}


\begin{solution}
There are 26 choices for each letter and 10 choices for each digit.

\begin{enumerate}[label=(\alph*)]
    \item With repetition allowed,
    \[
    26^2\,10^5=67{,}600{,}000.
    \]

    \item Without repetition, the two letter positions can be filled in
    \(26\cdot25\) ways and the five digit positions in
    \(10\cdot9\cdot8\cdot7\cdot6\) ways. Thus
    \[
    26\cdot25\cdot10\cdot9\cdot8\cdot7\cdot6
    =196{,}560{,}000.
    \]
\end{enumerate}
\end{solution}







   \newpage \begin{problem}
When all letters are used, how many different letter arrangements can be made from the letters in each word?
\begin{enumerate}[label=(\alph*)]
    \item FLUKE
    \item PROPOSE
    \item MISSISSIPPI
    \item ARRANGE
\end{enumerate}
\end{problem}

\begin{solution}
For \(n\) letters with repeated-letter multiplicities \(n_1,\ldots,n_r\), the number of distinct arrangements is
\[
\frac{n!}{n_1!\cdots n_r!}.
\]

\begin{enumerate}[label=(\alph*)]
    \item FLUKE has five distinct letters:
    \[
    5!=120.
    \]

    \item PROPOSE has seven letters, with \(P\) repeated twice and \(O\) repeated twice:
    \[
    \frac{7!}{2!\,2!}=1260.
    \]

    \item MISSISSIPPI has 11 letters, with
    \(I\) repeated four times, \(S\) four times, \(P\) twice, and \(M\) once:
    \[
    \frac{11!}{4!\,4!\,2!}=34{,}650.
    \]

    \item ARRANGE has seven letters, with \(A\) repeated twice and \(R\) repeated twice:
    \[
    \frac{7!}{2!\,2!}=1260.
    \]
\end{enumerate}
\end{solution}







    \newpage \begin{problem}
\begin{enumerate}[label=(\alph*)]
    \item Show that
    \[
    \sum_{k=0}^{n}\binom{n}{k}2^k=3^n.
    \]
    \item Simplify
    \[
    \sum_{k=0}^{n}\binom{n}{k}x^k.
    \]
\end{enumerate}
\end{problem}

\begin{solution}
The binomial theorem states
\[
(a+b)^n=\sum_{k=0}^{n}\binom{n}{k}a^{n-k}b^k.
\]

\begin{enumerate}[label=(\alph*)]
    \item Set \(a=1\) and \(b=2\):
    \[
    \sum_{k=0}^{n}\binom{n}{k}2^k
    =
    (1+2)^n
    =
    3^n.
    \]

    \item Set \(a=1\) and \(b=x\):
    \[
    \sum_{k=0}^{n}\binom{n}{k}x^k
    =
    (1+x)^n.
    \]
\end{enumerate}
\end{solution}
 





   \newpage \begin{problem}
If 12 people are to be divided into three committees of respective sizes \(3\), \(4\), and \(5\), how many divisions are possible?
\end{problem}
\begin{solution}
Choose 3 of the 12 people for the first committee, then 4 of the remaining 9 for the second; the final 5 form the third:
\[
\binom{12}{3}\binom{9}{4}\binom{5}{5}
=
\frac{12!}{3!\,4!\,5!}
=
27{,}720.
\]
\end{solution}


 






   \newpage \begin{problem}


\begin{enumerate}
    \item Expand $(3x^2+y)^5$. 
\item Without considering part (a), i.e., without using the expansion, write down the third term.
\end{enumerate}

\end{problem}


\begin{solution}
(1) By the binomial theorem,
\[
(3x^2+y)^5
=
\sum_{k=0}^{5}\binom{5}{k}(3x^2)^{5-k}y^k.
\]
Therefore,
\[
 {
(3x^2+y)^5
=
243x^{10}
+405x^8y
+270x^6y^2
+90x^4y^3
+15x^2y^4
+y^5.
}
\]


(2) The $(k+1)^{th}$ term is given by $\binom{n}{k} a^{k}b^{n-k}$. Thus, the third term is given by $\binom{5}{2} (y)^2(3x^2)^{5-2}.$
\end{solution}









   \newpage \begin{problem}
Prove that
\[
\binom{n+m}{r}
=
\binom{n}{0}\binom{m}{r}
+\binom{n}{1}\binom{m}{r-1}
+\cdots+
\binom{n}{r}\binom{m}{0}.
\]
\emph{Hint:} Consider a group of \(n\) men and \(m\) women and count groups of size \(r\).
\end{problem}
\begin{solution}
Count the number of ways to form a group of \(r\) people from \(n\) men and \(m\) women.

Directly, there are
\[
\binom{n+m}{r}
\]
such groups.

Alternatively, classify the group according to the number \(k\) of men selected. If exactly \(k\) men are selected, then \(r-k\) women are selected, giving
\[
\binom{n}{k}\binom{m}{r-k}
\]
possibilities. Summing over all feasible values of \(k\) gives
\[
\binom{n+m}{r}
=
\sum_{k=0}^{r}\binom{n}{k}\binom{m}{r-k}.
\]
Terms with impossible choices are interpreted as zero.
\end{solution}
 


   \newpage \begin{problem}
\begin{enumerate}[label=(\alph*)]
    \item Give a combinatorial proof of
    \[
    \sum_{k=1}^{n}k\binom{n}{k}=n2^{n-1}.
    \]
    Interpret both sides as the number of ways to select a nonempty committee from \(n\) people and choose its chairperson.

    \item Verify for \(n=1,2,3,4,5\), and then prove combinatorially, that
    \[
    \sum_{k=1}^{n}\binom{n}{k}k^2
    =
    2^{n-2}n(n+1).
    \]
    Interpret both sides as the number of ways to select a committee, its chairperson, and its secretary, where the chairperson and secretary may be the same person.

    \item Prove that
    \[
    \sum_{k=1}^{n}\binom{n}{k}k^3
    =
    2^{n-3}n^2(n+3).
    \]
\end{enumerate}
\end{problem}
\begin{solution}
\begin{enumerate}[label=(\alph*)]
    \item Count pairs consisting of a nonempty committee and its chairperson.

    If the committee has size \(k\), choose its members in \(\binom nk\) ways and its chairperson in \(k\) ways. This gives
    \[
    \sum_{k=1}^{n}k\binom nk.
    \]

    Alternatively, first choose the chairperson in \(n\) ways. Each of the other \(n-1\) people may independently be included or excluded, giving \(2^{n-1}\) choices. Hence
    \[
    \sum_{k=1}^{n}k\binom nk=n2^{n-1}.
    \]
Count selections of a committee, a chairperson, and a secretary, where the two officers may coincide. A committee of size \(k\) contributes \(k^2\binom nk\), so the total is the left-hand side.

    Split into two cases.

    If the chairperson and secretary are the same, choose that person in \(n\) ways and choose any subset of the remaining \(n-1\) people:
    \[
    n2^{n-1}.
    \]

    If they are different, choose the chairperson and secretary as an ordered pair in \(n(n-1)\) ways, and choose any subset of the remaining \(n-2\) people:
    \[
    n(n-1)2^{n-2}.
    \]

    Therefore,
    \[
    \sum_{k=1}^{n}\binom nk k^2
    =
    n2^{n-1}+n(n-1)2^{n-2}
    =
    2^{n-2}n(n+1).
    \]

    \item Count selections of a committee together with three ordered offices, where a person may hold more than one office. A committee of size \(k\) contributes \(k^3\binom nk\).

    Classify according to the number of distinct officeholders.

    One distinct officeholder:
    \[
    n2^{n-1}.
    \]

    Exactly two distinct officeholders: choose which one of the three offices is held by the person who holds only one office (\(3\) choices), then choose the two distinct people in order (\(n(n-1)\) choices), and choose any subset of the remaining \(n-2\) people:
    \[
    3n(n-1)2^{n-2}.
    \]

    Three distinct officeholders:
    \[
    n(n-1)(n-2)2^{n-3}.
    \]

    Adding,
    \begin{align*}
    \sum_{k=1}^{n}\binom nk k^3
    &=
    n2^{n-1}
    +3n(n-1)2^{n-2}
    +n(n-1)(n-2)2^{n-3}\\
    &=
    2^{n-3}n^2(n+3).
    \end{align*}
\end{enumerate}
\end{solution}
 


   \newpage \begin{problem}
Show that, for \(n>0\),
\[
\sum_{i=0}^{n}(-1)^i\binom{n}{i}=0.
\]
\end{problem}
\begin{solution}
Apply the binomial theorem to \((1-1)^n\):
\[
(1-1)^n
=
\sum_{i=0}^{n}\binom ni1^{n-i}(-1)^i
=
\sum_{i=0}^{n}(-1)^i\binom ni.
\]
For \(n>0\), the left-hand side is \(0^n=0\). Therefore,
\[
\sum_{i=0}^{n}(-1)^i\binom ni=0.
\]
\end{solution}
 


   \newpage \begin{problem}
A total of \(28\%\) of American males smoke cigarettes, \(7\%\) smoke cigars, and \(5\%\) smoke both cigars and cigarettes.
\begin{enumerate}[label=(\alph*)]
    \item What percentage smoke neither cigars nor cigarettes?
    \item What percentage smoke cigars but not cigarettes?
\end{enumerate}
\end{problem}
\begin{solution}
Let \(C\) be the event of smoking cigarettes and \(G\) the event of smoking cigars. Then
\[
P(C)=0.28,\qquad P(G)=0.07,\qquad P(C\cap G)=0.05.
\]

\begin{enumerate}[label=(\alph*)]
    \item
    \[
    P(C\cup G)=0.28+0.07-0.05=0.30.
    \]
    Thus
    \[
    P((C\cup G)^c)=1-0.30=0.70.
    \]
    The answer is \( {70\%}\).

    \item
    \[
    P(G\cap C^c)=P(G)-P(G\cap C)=0.07-0.05=0.02.
    \]
    The answer is \( {2\%}\).
\end{enumerate}
\end{solution}
 


   \newpage \begin{problem}
Two cards are chosen uniformly at random without replacement from a standard deck of 52 playing cards. What is the probability that they
\begin{enumerate}[label=(\alph*)]
    \item are both aces?
    \item have the same rank?
\end{enumerate}
\end{problem}
\begin{solution}
There are \(\binom{52}{2}\) unordered pairs of cards.

\begin{enumerate}[label=(\alph*)]
    \item There are \(\binom42\) pairs of aces:
    \[
    P(\text{both aces})
    =
    \frac{\binom42}{\binom{52}{2}}
    =
    \frac{6}{1326}
    =
     {\frac1{221}}.
    \]

    \item For each of the 13 ranks, choose 2 of the 4 suits:
    \[
    P(\text{same rank})
    =
    \frac{13\binom42}{\binom{52}{2}}
    =
    \frac{78}{1326}
    =
     {\frac1{17}}.
    \]
\end{enumerate}
\end{solution}
 

   \newpage \begin{problem}
Let \(E,F,\) and \(G\) be three events. Write expressions for the events that
\begin{enumerate}[label=(\alph*)]
    \item only \(E\) occurs;
    \item both \(E\) and \(G\), but not \(F\), occur;
    \item at least one of the events occurs;
    \item at least two of the events occur;
    \item all three events occur;
    \item none of the events occurs;
    \item at most one of the events occurs;
    \item at most two of the events occur;
    \item exactly two of the events occur;
    \item at most three of the events occur.
\end{enumerate}
\end{problem}
\begin{solution}
\begin{enumerate}[label=(\alph*)]
    \item Only \(E\): $E\cap F^c\cap G^c.$


    \item \(E\) and \(G\), but not \(F\): $    E\cap F^c\cap G.$
    

    \item At least one: $    E\cup F\cup G.
$ 

    \item At least two: $    (E\cap F)\cup(E\cap G)\cup(F\cap G).
$
 

    \item All three: $  E\cap F\cap G.$
  

    \item None: $    E^c\cap F^c\cap G^c=(E\cup F\cup G)^c.$
   

    \item At most one:
    \[
    (E^c\cap F^c\cap G^c)
    \cup(E\cap F^c\cap G^c)
    \cup(E^c\cap F\cap G^c)
    \cup(E^c\cap F^c\cap G).
    \]
    Equivalently,
    \[
    \bigl[(E\cap F)\cup(E\cap G)\cup(F\cap G)\bigr]^c.
    \]

    \item At most two:
    \[
    (E\cap F\cap G)^c.
    \]

    \item Exactly two:
    \[
    (E\cap F\cap G^c)
    \cup(E\cap F^c\cap G)
    \cup(E^c\cap F\cap G).
    \]

    \item At most three:
    \[
    \Omega,
    \]
    because there are only three events.
\end{enumerate}
\end{solution}
 

   \newpage \begin{problem}
If \(P(E)=0.9\) and \(P(F)=0.8\), show that \(P(E\cap F)\geq 0.7\). More generally, prove
\[
P(E\cap F)\geq P(E)+P(F)-1.
\]
\end{problem}
\begin{solution}
The addition rule gives
\[
P(E\cup F)=P(E)+P(F)-P(E\cap F).
\]
Because \(P(E\cup F)\leq1\),
\[
P(E)+P(F)-P(E\cap F)\leq1,
\]
and therefore
\[
P(E\cap F)\geq P(E)+P(F)-1.
\]
For \(P(E)=0.9\) and \(P(F)=0.8\),
\[
P(E\cap F)\geq0.9+0.8-1=0.7.
\]
\end{solution}
 

   \newpage \begin{problem}
Prove that
\[
P(E\cap F^c)=P(E)-P(E\cap F).
\]
\end{problem}
\begin{solution}
The event \(E\) is the disjoint union
\[
E=(E\cap F)\cup(E\cap F^c).
\]
Hence, by additivity,
\[
P(E)=P(E\cap F)+P(E\cap F^c).
\]
Rearranging gives
\[
P(E\cap F^c)=P(E)-P(E\cap F).
\]
\end{solution}
 



   \newpage \begin{problem}
A laboratory blood test has sensitivity \(0.95\): when a person has a certain disease, the test is positive with probability \(0.95\). The false-positive probability is \(0.01\): when a person is healthy, the test is positive with probability \(0.01\). If \(0.5\%\) of the population has the disease, find the probability that a person has the disease given that the test result is positive.
\end{problem}
\begin{solution}
Let \(D\) be the event that the person has the disease and \(+\) the event that the test is positive. Then
\[
P(D)=0.005,\quad
P(D^c)=0.995,\quad
P(+\mid D)=0.95,\quad
P(+\mid D^c)=0.01.
\]
By Bayes' rule,
\begin{align*}
P(D\mid +)
&=
\frac{P(+\mid D)P(D)}
{P(+\mid D)P(D)+P(+\mid D^c)P(D^c)}\\
&=
\frac{(0.95)(0.005)}
{(0.95)(0.005)+(0.01)(0.995)}\\
&=
\frac{0.00475}{0.01470}
=
\frac{95}{294}
\approx
 {0.3231}.
\end{align*}
Thus, given a positive result, the probability of having the disease is about \(32.3\%\).
\end{solution}
 


   \newpage \begin{problem}
Let \(A\subseteq B\). Express the following probabilities as simply as possible:
\[
P(A\mid B),\qquad
P(A\mid B^c),\qquad
P(B\mid A),\qquad
P(B\mid A^c).
\]
State any conditions needed for the conditional probabilities to be defined.
\end{problem}
\begin{solution}
Because \(A\subseteq B\), we have \(A\cap B=A\), \(A\cap B^c=\varnothing\), and \(B\cap A=A\).

Provided \(P(B)>0\),
\[
P(A\mid B)
=
\frac{P(A\cap B)}{P(B)}
=
 {\frac{P(A)}{P(B)}}.
\]

Provided \(P(B^c)>0\),
\[
P(A\mid B^c)
=
\frac{P(A\cap B^c)}{P(B^c)}
=
 {0}.
\]

Provided \(P(A)>0\),
\[
P(B\mid A)
=
\frac{P(B\cap A)}{P(A)}
=
 {1}.
\]

Provided \(P(A^c)>0\),
\[
P(B\mid A^c)
=
\frac{P(B\cap A^c)}{P(A^c)}
=
 {\frac{P(B)-P(A)}{1-P(A)}}.
\]
\end{solution}
 

   \newpage \begin{problem}
Two fair dice are rolled, and \(X\) is the product of the two outcomes. Compute \(P(X=i)\) for \(i=1,\ldots,36\).
\end{problem}
\begin{solution}
The 36 ordered outcomes are equally likely. The nonzero probabilities are
\[
\begin{array}{c|cccccccccccccccccc}
i
&1&2&3&4&5&6&8&9&10&12&15&16&18&20&24&25&30&36\\ \hline
36P(X=i)
&1&2&2&3&2&4&2&1&2&4&2&1&2&2&2&1&2&1
\end{array}
\]
Thus
\[
P(X=i)=
\begin{cases}
\dfrac{1}{36},
&i\in\{1,9,16,25,36\},\\[1mm]
\dfrac{1}{18},
&i\in\{2,3,5,8,10,15,18,20,24,30\},\\[1mm]
\dfrac{1}{12},
&i=4,\\[1mm]
\dfrac{1}{9},
&i\in\{6,12\},\\[1mm]
0,
&\text{otherwise}.
\end{cases}
\]
\end{solution}
 

   \newpage \begin{problem}
Suppose the distribution function of \(X\) is
\[
F(b)=
\begin{cases}
0, & b<0,\\[2mm]
\dfrac{b}{4}, & 0\leq b<1,\\[2mm]
\dfrac12+\dfrac{b-1}{4}, & 1\leq b<2,\\[2mm]
\dfrac{11}{12}, & 2\leq b<3,\\[2mm]
1, & b\geq 3.
\end{cases}
\]
\begin{enumerate}[label=(\alph*)]
    \item Find \(P(X=i)\) for \(i=1,2,3\).
    \item Find
    \[
    P\left(\frac12<X<\frac32\right).
    \]
\end{enumerate}
\end{problem}
\begin{solution}
For any point \(a\),
\[
P(X=a)=F(a)-F(a^-),
\]
the size of the jump of the CDF at \(a\).

\begin{enumerate}[label=(\alph*)]
    \item At \(1\),
    \[
    P(X=1)
    =
    F(1)-F(1^-)
    =
    \frac12-\frac14
    =
     {\frac14}.
    \]

    At \(2\),
    \[
    P(X=2)
    =
    F(2)-F(2^-)
    =
    \frac{11}{12}-\frac34
    =
     {\frac16}.
    \]

    At \(3\),
    \[
    P(X=3)
    =
    F(3)-F(3^-)
    =
    1-\frac{11}{12}
    =
     {\frac1{12}}.
    \]

    \item Since there is no point mass at \(1/2\) or \(3/2\),
    \begin{align*}
    P\left(\frac12<X<\frac32\right)
    &=
    F\left(\frac32\right)-F\left(\frac12\right)\\
    &=
    \left(\frac12+\frac{(3/2)-1}{4}\right)
    -\frac{1/2}{4}\\
    &=
    \frac58-\frac18
    =
     {\frac12}.
    \end{align*}
\end{enumerate}
\end{solution}
 

   \newpage \begin{problem}
Suppose
\[
F(b)=
\begin{cases}
0, & b<0,\\[1mm]
\dfrac12, & 0\leq b<1,\\[1mm]
\dfrac35, & 1\leq b<2,\\[1mm]
\dfrac45, & 2\leq b<3,\\[1mm]
\dfrac9{10}, & 3\leq b<3.5,\\[1mm]
1, & b\geq 3.5.
\end{cases}
\]
Calculate the probability mass function of \(X\).
\end{problem}
\begin{solution}
The probability at each support point is the jump of the CDF:
\[
p_X(x)=P(X=x)=F(x)-F(x^-).
\]
Therefore,
\[
P(X=0)=\frac12,
\]
\[
P(X=1)=\frac35-\frac12=\frac1{10},
\]
\[
P(X=2)=\frac45-\frac35=\frac15,
\]
\[
P(X=3)=\frac9{10}-\frac45=\frac1{10},
\]
and
\[
P(X=3.5)=1-\frac9{10}=\frac1{10}.
\]
Hence
\[
 {
p_X(x)=
\begin{cases}
\dfrac12, & x=0,\\[1mm]
\dfrac1{10}, & x=1,\\[1mm]
\dfrac15, & x=2,\\[1mm]
\dfrac1{10}, & x=3,\\[1mm]
\dfrac1{10}, & x=3.5,\\[1mm]
0, & \text{otherwise}.
\end{cases}}
\]
The probabilities sum to
\[
\frac12+\frac1{10}+\frac15+\frac1{10}+\frac1{10}=1.
\]
\end{solution}
 


%%%%%%%%%%%%%%%%%%%%%%%%%%%%%%%%%%%%%%%%%%%%%%%%%%%%%%%%%%%%%%
%%%%%%%%%%%%%%%%%%%%%%  PROBLEM 24  %%%%%%%%%%%%%%%%%%%%%%%%%%
%%%%%%%%%%%%%%%%%%%%%%%%%%%%%%%%%%%%%%%%%%%%%%%%%%%%%%%%%%%%%%

   \newpage \begin{problem}

The expected number of typographical errors on a page of a certain
magazine is \(0.2\). Assuming that the number of errors on a page has a
Poisson distribution, find the probability that the next page contains

\begin{enumerate}[label=(\alph*)]
    \item no typographical errors;
    \item two or more typographical errors.
\end{enumerate}
\end{problem}

\begin{solution}

Let \(X\) denote the number of typographical errors on the next page.
Under the Poisson assumption,

\[
X\sim\operatorname{Pois}(0.2).
\]

Thus,

\[
P(X=k)=e^{-0.2}\frac{(0.2)^k}{k!}.
\]

\begin{enumerate}[label=(\alph*)]

\item

\[
P(X=0)
=
e^{-0.2}
\approx
 {0.8187}.
\]

\item Using the complement,

\[
\begin{aligned}
P(X\geq 2)
&=
1-P(X=0)-P(X=1)\\
&=
1-e^{-0.2}-(0.2)e^{-0.2}\\
&=
1-1.2e^{-0.2}\\
&\approx
 {0.0175}.
\end{aligned}
\]

\end{enumerate}

The expected value alone does not uniquely determine these probabilities.
The calculation requires the additional assumption that the number of
errors follows a Poisson distribution.

\end{solution}
 


%%%%%%%%%%%%%%%%%%%%%%%%%%%%%%%%%%%%%%%%%%%%%%%%%%%%%%%%%%%%%%
%%%%%%%%%%%%%%%%%%%%%%  PROBLEM 25  %%%%%%%%%%%%%%%%%%%%%%%%%%
%%%%%%%%%%%%%%%%%%%%%%%%%%%%%%%%%%%%%%%%%%%%%%%%%%%%%%%%%%%%%%

   \newpage \begin{problem}

The monthly worldwide average number of airplane crashes involving
commercial airlines is \(3.5\). Assuming a Poisson model, find the
probability that there will be

\begin{enumerate}[label=(\alph*)]
    \item at least two such accidents in the next month;
    \item at most one such accident in the next month.
\end{enumerate}
\end{problem}

\begin{solution}

Let \(X\) denote the number of accidents during the next month. Under the
Poisson assumption,

\[
X\sim\operatorname{Pois}(3.5).
\]

Therefore,

\[
P(X=k)=e^{-3.5}\frac{3.5^k}{k!}.
\]

\begin{enumerate}[label=(\alph*)]

\item

\[
\begin{aligned}
P(X\geq 2)
&=
1-P(X=0)-P(X=1)\\
&=
1-e^{-3.5}-3.5e^{-3.5}\\
&=
1-4.5e^{-3.5}\\
&\approx
 {0.8641}.
\end{aligned}
\]

\item

\[
\begin{aligned}
P(X\leq 1)
&=
P(X=0)+P(X=1)\\
&=
e^{-3.5}+3.5e^{-3.5}\\
&=
4.5e^{-3.5}\\
&\approx
 {0.1359}.
\end{aligned}
\]

Notice that the answers in parts (a) and (b) are complements.

\end{enumerate}

\end{solution}
 


%%%%%%%%%%%%%%%%%%%%%%%%%%%%%%%%%%%%%%%%%%%%%%%%%%%%%%%%%%%%%%
%%%%%%%%%%%%%%%%%%%%%%  PROBLEM 26  %%%%%%%%%%%%%%%%%%%%%%%%%%
%%%%%%%%%%%%%%%%%%%%%%%%%%%%%%%%%%%%%%%%%%%%%%%%%%%%%%%%%%%%%%

   \newpage \begin{problem}

Approximately \(80{,}000\) marriages took place in the state of New York
last year. Estimate the probability that, for at least one of these
couples,

\begin{enumerate}[label=(\alph*)]
    \item both partners were born on April 30;
    \item both partners celebrate their birthdays on the same day of the
    year.
\end{enumerate}

State the assumptions used.
\end{problem}
\begin{solution}

Assume that:

\begin{itemize}
    \item birthdays are uniformly distributed among the \(365\) days of
    the year;
    \item leap days are ignored;
    \item the birthdays of the two partners are independent;
    \item different couples may be treated independently.
\end{itemize}

Let \(n=80{,}000\).

\begin{enumerate}[label=(\alph*)]

\item For a particular couple, the probability that both partners were
born on April 30 is

\[
p=\left(\frac{1}{365}\right)^2.
\]

Thus, the exact probability that at least one couple has this property is

\[
\begin{aligned}
P(\text{at least one})
&=
1-(1-p)^{80{,}000}\\
&=
1-\left(1-\frac{1}{365^2}\right)^{80{,}000}.
\end{aligned}
\]

Because \(n\) is large and \(p\) is small, a Poisson approximation with

\[
\lambda=np
=
\frac{80{,}000}{365^2}
\approx0.60049
\]

gives

\[
P(\text{at least one})
\approx
1-e^{-0.60049}
\approx
 {0.4515}.
\]

\item For a particular couple, condition on the birthday of the first
partner. The probability that the second partner has the same birthday is

\[
p=\frac{1}{365}.
\]

Therefore,

\[
P(\text{at least one matching couple})
=
1-\left(1-\frac{1}{365}\right)^{80{,}000}.
\]

Using a Poisson approximation,

\[
\lambda
=
\frac{80{,}000}{365}
\approx219.18,
\]

so

\[
P(\text{at least one matching couple})
\approx
1-e^{-219.18}\approx 1.
\]

 


\end{enumerate}

\end{solution}
 


%%%%%%%%%%%%%%%%%%%%%%%%%%%%%%%%%%%%%%%%%%%%%%%%%%%%%%%%%%%%%%
%%%%%%%%%%%%%%%%%%%%%%  PROBLEM 27  %%%%%%%%%%%%%%%%%%%%%%%%%%
%%%%%%%%%%%%%%%%%%%%%%%%%%%%%%%%%%%%%%%%%%%%%%%%%%%%%%%%%%%%%%

   \newpage \begin{problem}

Compare the Poisson approximation with the exact binomial probability in
each of the following cases:

\begin{enumerate}[label=(\alph*)]
    \item \(P(X=2)\), where \(n=8\) and \(p=0.1\);
    \item \(P(X=9)\), where \(n=10\) and \(p=0.95\);
    \item \(P(X=0)\), where \(n=10\) and \(p=0.1\);
    \item \(P(X=4)\), where \(n=9\) and \(p=0.2\).
\end{enumerate}
\end{problem}

\begin{solution}

If $X\sim\operatorname{Bin}(n,p),
$ then the exact probability is

\[
P(X=k)=\binom nkp^k(1-p)^{n-k}.
\]

The usual Poisson approximation uses

\[
Y\sim\operatorname{Pois}(\lambda),
\qquad
\lambda=np,
\]

and

\[
P(X=k)\approx P(Y=k)
=
e^{-\lambda}\frac{\lambda^k}{k!}.
\]

\begin{enumerate}[label=(\alph*)]

\item Here \(n=8\), \(p=0.1\), \(k=2\), and \(\lambda=0.8\). The exact probability is

\[
\begin{aligned}
P(X=2)
&=
\binom82(0.1)^2(0.9)^6\\
&\approx
 {0.1488}.
\end{aligned}
\]
The Poisson approximation is reasonably accurate:
\[
\begin{aligned}
P(Y=2)
&=
e^{-0.8}\frac{0.8^2}{2!}\\
&\approx
 {0.1438}.
\end{aligned}
\]

\item Here \(n=10\), \(p=0.95\), \(k=9\), and \(\lambda=9.5\). The exact probability is

\[
\begin{aligned}
P(X=9)
&=
\binom{10}{9}(0.95)^9(0.05)\\
&\approx
 {0.3151}.
\end{aligned}
\]

The direct Poisson approximation is
poor because \(p=0.95\) is not small.
\[
\begin{aligned}
P(Y=9)
&=
e^{-9.5}\frac{9.5^9}{9!}\\
&\approx
 {0.1300}.
\end{aligned}
\]

 

A better approach is to let $Z=10-X,$ the number of failures. Then $Z\sim\operatorname{Bin}(10,0.05),
$ and \(X=9\) is equivalent to \(Z=1\). Approximating \(Z\) by
\(\operatorname{Pois}(0.5)\) gives the following, which is much closer to the exact value:
\[
P(X=9)
=
P(Z=1)
\approx
e^{-0.5}(0.5)
\approx0.3033,
\]

\item Here \(n=10\), \(p=0.1\), \(k=0\), and \(\lambda=1\). The exact probability is
\[
P(X=0)
=
(0.9)^{10}
\approx
 {0.3487}.
\]

The Poisson approximation is reasonably accurate: $P(Y=0)
=
e^{-1}
\approx
 {0.3679}.$
 

\item Here \(n=9\), \(p=0.2\), \(k=4\), and \(\lambda=1.8\). The exact probability is

\[
\begin{aligned}
P(X=4)
&=
\binom94(0.2)^4(0.8)^5\\
&\approx
 {0.0661}.
\end{aligned}
\]

The Poisson approximation is 

\[
\begin{aligned}
P(Y=4)
&=
e^{-1.8}\frac{1.8^4}{4!}\\
&\approx
 {0.0723}.
\end{aligned}
\]

\end{enumerate}

\end{solution}
 


%%%%%%%%%%%%%%%%%%%%%%%%%%%%%%%%%%%%%%%%%%%%%%%%%%%%%%%%%%%%%%
%%%%%%%%%%%%%%%%%%%%%%  PROBLEM 28  %%%%%%%%%%%%%%%%%%%%%%%%%%
%%%%%%%%%%%%%%%%%%%%%%%%%%%%%%%%%%%%%%%%%%%%%%%%%%%%%%%%%%%%%%

   \newpage \begin{problem}

The rate of a certain event in a state is one event per \(100{,}000\)
inhabitants per month.

\begin{enumerate}[label=(\alph*)]

    \item Find the probability that a city of \(400{,}000\) inhabitants
    experiences eight or more events in a given month.

    \item Find the probability that at least two months during a
    twelve-month year have eight or more events.

    \item Counting the present month as month \(1\), find the probability
    that the first month having eight or more events is month \(i\), where
    \(i\geq1\).

\end{enumerate}

State the assumptions used.
\end{problem}
\begin{solution}

Assume that monthly event counts are independent and identically
distributed Poisson random variables and that the rate remains constant
over time.

For a city of \(400{,}000\) inhabitants, the expected monthly number is

\[
\lambda
=
400{,}000\left(\frac{1}{100{,}000}\right)
=
4.
\]

Thus, if \(X\) is the number of events in a given month,

\[
X\sim\operatorname{Pois}(4).
\]

\begin{enumerate}[label=(\alph*)]

\item Let $q=P(X\geq8).$ Then

\[
\begin{aligned}
q
&=
1-P(X\leq7)\\
&=
1-\sum_{k=0}^{7}e^{-4}\frac{4^k}{k!}\\
&\approx
 {0.05113}.
\end{aligned}
\]

\item Let \(Y\) be the number of months, among the next 12 months, that
have at least eight events. Under the independence assumption,

\[
Y\sim\operatorname{Bin}(12,q),
\]

where \(q\approx0.05113\).

Therefore,

\[
\begin{aligned}
P(Y\geq2)
&=
1-P(Y=0)-P(Y=1)\\
&=
1-(1-q)^{12}
-12q(1-q)^{11}\\
&\approx
 {0.1229}.
\end{aligned}
\]

\item Let \(T\) denote the first month with at least eight events. Each
month is a success with probability \(q\), independently of other months.
Therefore,

\[
T\sim\operatorname{Geom}(q),
\]

using the convention that \(T\) is the trial number of the first success.

Hence,

\[
 {
P(T=i)=(1-q)^{i-1}q,
\qquad i=1,2,\ldots,
}
\]

where

\[
q
=
1-\sum_{k=0}^{7}e^{-4}\frac{4^k}{k!}
\approx0.05113.
\]

Numerically, $P(T=i)
\approx
(0.94887)^{i-1}(0.05113).$

 

\end{enumerate}

\end{solution}





%%%%%%%%%%%%%%%%%%%%%%%%%%%%%%%%%%%%%%%%%%%%%%%%%%%%%%%%%%%%%%
%%%%%%%%%%%%%%%%%%%%%%  PROBLEM 24  %%%%%%%%%%%%%%%%%%%%%%%%%%
%%%%%%%%%%%%%%%%%%%%%%%%%%%%%%%%%%%%%%%%%%%%%%%%%%%%%%%%%%%%%%

   \newpage \begin{problem}

The expected number of typographical errors on a page of a certain
magazine is \(0.2\). Assuming that the number of errors on a page has a
Poisson distribution, find the probability that the next page contains

\begin{enumerate}[label=(\alph*)]
    \item no typographical errors;
    \item two or more typographical errors.
\end{enumerate}
\end{problem}
\begin{solution}

Let \(X\) denote the number of typographical errors on the next page.
Under the Poisson assumption,

\[
X\sim\operatorname{Pois}(0.2).
\]

Thus,

\[
P(X=k)=e^{-0.2}\frac{(0.2)^k}{k!}.
\]

\begin{enumerate}[label=(\alph*)]

\item

\[
P(X=0)
=
e^{-0.2}
\approx
 {0.8187}.
\]

\item Using the complement,

\[
\begin{aligned}
P(X\geq 2)
&=
1-P(X=0)-P(X=1)\\
&=
1-e^{-0.2}-(0.2)e^{-0.2}\\
&=
1-1.2e^{-0.2}\\
&\approx
 {0.0175}.
\end{aligned}
\]

\end{enumerate}

The expected value alone does not uniquely determine these probabilities.
The calculation requires the additional assumption that the number of
errors follows a Poisson distribution.

\end{solution}
 


%%%%%%%%%%%%%%%%%%%%%%%%%%%%%%%%%%%%%%%%%%%%%%%%%%%%%%%%%%%%%%
%%%%%%%%%%%%%%%%%%%%%%  PROBLEM 25  %%%%%%%%%%%%%%%%%%%%%%%%%%
%%%%%%%%%%%%%%%%%%%%%%%%%%%%%%%%%%%%%%%%%%%%%%%%%%%%%%%%%%%%%%



%%%%%%%%%%%%%%%%%%%%%%%%%%%%%%%%%%%%%%%%%%%%%%%%%%%%%%%%%%%%%%
%%%%%%%%%%%%%%%%%%%%%%  PROBLEM 26  %%%%%%%%%%%%%%%%%%%%%%%%%%
%%%%%%%%%%%%%%%%%%%%%%%%%%%%%%%%%%%%%%%%%%%%%%%%%%%%%%%%%%%%%%



%%%%%%%%%%%%%%%%%%%%%%%%%%%%%%%%%%%%%%%%%%%%%%%%%%%%%%%%%%%%%%
%%%%%%%%%%%%%%%%%%%%%%  PROBLEM 27  %%%%%%%%%%%%%%%%%%%%%%%%%%
%%%%%%%%%%%%%%%%%%%%%%%%%%%%%%%%%%%%%%%%%%%%%%%%%%%%%%%%%%%%%%



%%%%%%%%%%%%%%%%%%%%%%%%%%%%%%%%%%%%%%%%%%%%%%%%%%%%%%%%%%%%%%
%%%%%%%%%%%%%%%%%%%%%%  PROBLEM 28  %%%%%%%%%%%%%%%%%%%%%%%%%%
%%%%%%%%%%%%%%%%%%%%%%%%%%%%%%%%%%%%%%%%%%%%%%%%%%%%%%%%%%%%%%






%%%%%%%%%%%%%%%%%%%%%%%%%%%%%%%%%%%%%%%%%%%%%%%%%%%%%%%%%%%%%%
%%%%%%%%%%%%%%%%%%%%%%  PROBLEM 29  %%%%%%%%%%%%%%%%%%%%%%%%%%
%%%%%%%%%%%%%%%%%%%%%%%%%%%%%%%%%%%%%%%%%%%%%%%%%%%%%%%%%%%%%%

   \newpage \begin{problem}

A random variable \(X\) is said to have the Yule--Simon distribution if

\[
P(X=n)=\frac{4}{n(n+1)(n+2)},
\qquad n=1,2,\ldots
\]

\begin{enumerate}[label=(\alph*)]
    \item Show that this is a probability mass function.
    \item Show that \(E[X]=2\).
    \item Show that \(E[X^2]=\infty\).
\end{enumerate}

 \hrule 

The proposed probabilities are nonnegative. It remains to verify that
they sum to one.

\begin{enumerate}[label=(\alph*)]

\item First observe that

\[
\frac{1}{n(n+1)(n+2)}
=
\frac12
\left(
\frac{1}{n(n+1)}
-
\frac{1}{(n+1)(n+2)}
\right).
\]

Therefore,

\[
\frac{4}{n(n+1)(n+2)}
=
2
\left(
\frac{1}{n(n+1)}
-
\frac{1}{(n+1)(n+2)}
\right).
\]

For \(N\geq1\),

\begin{align*}
\sum_{n=1}^{N}\frac{4}{n(n+1)(n+2)}
&=
2\sum_{n=1}^{N}
\left(
\frac{1}{n(n+1)}
-
\frac{1}{(n+1)(n+2)}
\right)\\
&=
2\left(
\frac{1}{1\cdot2}
-
\frac{1}{(N+1)(N+2)}
\right).
\end{align*}

Letting \(N\to\infty\),

\[
\sum_{n=1}^{\infty}\frac{4}{n(n+1)(n+2)}
=
2\left(\frac12\right)
=
1.
\]

Hence the given function is a probability mass function.

\item

\begin{align*}
E[X]
&=
\sum_{n=1}^{\infty}
n\frac{4}{n(n+1)(n+2)}\\
&=
4\sum_{n=1}^{\infty}
\frac{1}{(n+1)(n+2)}.
\end{align*}

Using

\[
\frac{1}{(n+1)(n+2)}
=
\frac{1}{n+1}-\frac{1}{n+2},
\]

we obtain

\begin{align*}
E[X]
&=
4\sum_{n=1}^{\infty}
\left(
\frac{1}{n+1}-\frac{1}{n+2}
\right)\\
&=
4\left(\frac12\right)\\
&=
 {2}.
\end{align*}

\item

\begin{align*}
E[X^2]
&=
\sum_{n=1}^{\infty}
n^2\frac{4}{n(n+1)(n+2)}\\
&=
4\sum_{n=1}^{\infty}
\frac{n}{(n+1)(n+2)}.
\end{align*}

For \(n\geq2\),

\[
n+1\leq2n,
\qquad
n+2\leq2n,
\]

and hence

\[
\frac{n}{(n+1)(n+2)}
\geq
\frac{n}{(2n)(2n)}
=
\frac{1}{4n}.
\]

Therefore,

\[
E[X^2]
\geq
4\sum_{n=2}^{\infty}\frac{1}{4n}
=
\sum_{n=2}^{\infty}\frac1n
=
\infty.
\]

Thus,

\[
 {E[X^2]=\infty}.
\]

In particular, \(X\) has a finite mean but does not have a finite
variance.

\end{enumerate}

 
\end{problem}


%%%%%%%%%%%%%%%%%%%%%%%%%%%%%%%%%%%%%%%%%%%%%%%%%%%%%%%%%%%%%%
%%%%%%%%%%%%%%%%%%%%%%  PROBLEM 30  %%%%%%%%%%%%%%%%%%%%%%%%%%
%%%%%%%%%%%%%%%%%%%%%%%%%%%%%%%%%%%%%%%%%%%%%%%%%%%%%%%%%%%%%%

   \newpage \begin{problem}

Let \(X\) be a random variable with

\[
E[X]=\mu
\qquad\text{and}\qquad
\operatorname{Var}(X)=\sigma^2,
\]

where \(\sigma>0\). Define

\[
Y=\frac{X-\mu}{\sigma}.
\]

Find \(E[Y]\) and \(\operatorname{Var}(Y)\).
\end{problem}
\begin{solution}

Using linearity of expectation,

\begin{align*}
E[Y]
&=
E\left[\frac{X-\mu}{\sigma}\right]\\
&=
\frac{1}{\sigma}
\left(E[X]-\mu\right)\\
&=
\frac{1}{\sigma}(\mu-\mu)\\
&=
 {0}.
\end{align*}

For the variance, use

\[
\operatorname{Var}(aX+b)
=
a^2\operatorname{Var}(X).
\]

Thus,

\begin{align*}
\operatorname{Var}(Y)
&=
\operatorname{Var}
\left(
\frac{X-\mu}{\sigma}
\right)\\
&=
\frac{1}{\sigma^2}\operatorname{Var}(X)\\
&=
\frac{\sigma^2}{\sigma^2}\\
&=
 {1}.
\end{align*}

Therefore, the standardized random variable \(Y\) has mean \(0\) and
variance \(1\).

\end{solution}
 


%%%%%%%%%%%%%%%%%%%%%%%%%%%%%%%%%%%%%%%%%%%%%%%%%%%%%%%%%%%%%%
%%%%%%%%%%%%%%%%%%%%%%  PROBLEM 31  %%%%%%%%%%%%%%%%%%%%%%%%%%
%%%%%%%%%%%%%%%%%%%%%%%%%%%%%%%%%%%%%%%%%%%%%%%%%%%%%%%%%%%%%%

   \newpage \begin{problem}

Show how the binomial probability formula

\[
P(X=i)
=
\binom ni p^i(1-p)^{n-i},
\qquad i=0,\ldots,n,
\]

leads to a proof of the binomial theorem

\[
(x+y)^n
=
\sum_{i=0}^{n}
\binom ni x^iy^{n-i}
\]

when \(x\) and \(y\) are nonnegative.
\end{problem}
\begin{solution}

First suppose that \(x+y>0\), and define

\[
p=\frac{x}{x+y}.
\]

Then

\[
1-p=\frac{y}{x+y}.
\]

If \(X\sim\operatorname{Bin}(n,p)\), the probabilities of all possible
values of \(X\) sum to one:

\[
\sum_{i=0}^{n}
\binom ni p^i(1-p)^{n-i}
=
1.
\]

Substituting the expressions for \(p\) and \(1-p\),

\[
\sum_{i=0}^{n}
\binom ni
\left(\frac{x}{x+y}\right)^i
\left(\frac{y}{x+y}\right)^{n-i}
=
1.
\]

Because every term has denominator \((x+y)^n\),

\[
\frac{1}{(x+y)^n}
\sum_{i=0}^{n}
\binom ni x^iy^{n-i}
=
1.
\]

Multiplying by \((x+y)^n\) gives

\[
 {
(x+y)^n
=
\sum_{i=0}^{n}
\binom ni x^iy^{n-i}.
}
\]

If \(x=y=0\), the identity also holds directly for \(n>0\), because both
sides equal zero.

\end{solution}
 


%%%%%%%%%%%%%%%%%%%%%%%%%%%%%%%%%%%%%%%%%%%%%%%%%%%%%%%%%%%%%%
%%%%%%%%%%%%%%%%%%%%%%  PROBLEM 32  %%%%%%%%%%%%%%%%%%%%%%%%%%
%%%%%%%%%%%%%%%%%%%%%%%%%%%%%%%%%%%%%%%%%%%%%%%%%%%%%%%%%%%%%%

   \newpage \begin{problem}

Let \(X\) be a Poisson random variable with parameter \(\lambda>0\).
Show that \(P(X=i)\) first increases and then decreases as \(i\)
increases. Determine where the maximum occurs.
\end{problem}
\begin{solution}

The Poisson probability mass function is

\[
p_i=P(X=i)=e^{-\lambda}\frac{\lambda^i}{i!},
\qquad i=0,1,2,\ldots
\]

For \(i\geq1\), consider the ratio of consecutive probabilities:

\begin{align*}
\frac{p_i}{p_{i-1}}
&=
\frac{e^{-\lambda}\lambda^i/i!}
{e^{-\lambda}\lambda^{i-1}/(i-1)!}\\
&=
\frac{\lambda}{i}.
\end{align*}

Therefore,

\[
p_i>p_{i-1}
\quad\Longleftrightarrow\quad
i<\lambda,
\]

and

\[
p_i<p_{i-1}
\quad\Longleftrightarrow\quad
i>\lambda.
\]

Thus the probabilities increase while \(i<\lambda\) and decrease once
\(i>\lambda\).

If \(\lambda\) is not an integer, the unique mode is

\[
 {\lfloor\lambda\rfloor}.
\]

If \(\lambda=m\) is a positive integer, then

\[
\frac{p_m}{p_{m-1}}
=
\frac{m}{m}
=
1,
\]

so

\[
p_{m-1}=p_m.
\]

Hence there are two modes:

\[
 {m-1\ \text{and}\ m}.
\]

In particular, \(\lfloor\lambda\rfloor\) is always a mode, but when
\(\lambda\) is an integer, it is not the only mode.

\end{solution}
 


%%%%%%%%%%%%%%%%%%%%%%%%%%%%%%%%%%%%%%%%%%%%%%%%%%%%%%%%%%%%%%
%%%%%%%%%%%%%%%%%%%%%%  PROBLEM 33  %%%%%%%%%%%%%%%%%%%%%%%%%%
%%%%%%%%%%%%%%%%%%%%%%%%%%%%%%%%%%%%%%%%%%%%%%%%%%%%%%%%%%%%%%



%%%%%%%%%%%%%%%%%%%%%%%%%%%%%%%%%%%%%%%%%%%%%%%%%%%%%%%%%%%%%%
%%%%%%%%%%%%%%%%%%%%%%  PROBLEM 34  %%%%%%%%%%%%%%%%%%%%%%%%%%
%%%%%%%%%%%%%%%%%%%%%%%%%%%%%%%%%%%%%%%%%%%%%%%%%%%%%%%%%%%%%%

   \newpage \begin{problem}

Let \(X\) be a geometric random variable with success probability \(p\),
where \(X\) is the trial number on which the first success occurs. Show
analytically that

\[
P(X=n+k\mid X>n)=P(X=k),
\qquad n\geq0,\quad k\geq1.
\]

Also give a verbal explanation.
\end{problem}
\begin{solution}

Because \(X\) is geometric,

\[
P(X=j)=(1-p)^{j-1}p,
\qquad j=1,2,\ldots
\]

Also,

\[
P(X>n)=(1-p)^n.
\]

Since the event \(\{X=n+k\}\) is contained in the event \(\{X>n\}\),

\begin{align*}
P(X=n+k\mid X>n)
&=
\frac{P(X=n+k)}{P(X>n)}\\
&=
\frac{(1-p)^{n+k-1}p}{(1-p)^n}\\
&=
(1-p)^{k-1}p\\
&=
P(X=k).
\end{align*}

Therefore,

\[
 {
P(X=n+k\mid X>n)=P(X=k).
}
\]

For the verbal argument, the condition \(X>n\) means that the first
\(n\) trials were failures. Because the trials are independent, those
failures do not affect the outcomes of future trials. Starting after
trial \(n\), the waiting time until the first success has the same
geometric distribution as it had at the beginning. This property is
called the \emph{memoryless property}.

\end{solution}
 



 
%%%%%%%%%%%%%%%%%%%%%%%%%%%%%%%%%%%%%%%%%%%%%%%%%%%%%%%%%%%%%%
%%%%%%%%%%%%%%%%%%%%%%  PROBLEM 35  %%%%%%%%%%%%%%%%%%%%%%%%%%
%%%%%%%%%%%%%%%%%%%%%%%%%%%%%%%%%%%%%%%%%%%%%%%%%%%%%%%%%%%%%%

   \newpage \begin{problem}

Compute \(E[X]\) for each of the following density functions.

\begin{enumerate}[label=(\alph*)]

\item
\[
f(x)=
\begin{cases}
\dfrac14xe^{-x/2}, & x>0,\\[1mm]
0, & \text{otherwise}.
\end{cases}
\]

\item
\[
f(x)=
\begin{cases}
c(1-x^2), & -1<x<1,\\[1mm]
0, & \text{otherwise}.
\end{cases}
\]

\item
\[
f(x)=
\begin{cases}
\dfrac{5}{x^2}, & x>5,\\[1mm]
0, & x\leq5.
\end{cases}
\]

\end{enumerate}
\end{problem}
\begin{solution}

\begin{enumerate}[label=(\alph*)]

\item By definition,

\[
E[X]
=
\int_{0}^{\infty}x f(x)\,dx.
\]

Therefore,

\begin{align*}
E[X]
&=
\int_{0}^{\infty}
x\left(\frac14xe^{-x/2}\right)\,dx\\
&=
\frac14
\int_{0}^{\infty}x^2e^{-x/2}\,dx.
\end{align*}
Using integration by parts, we have

\[
\int_{0}^{\infty}x^2e^{-x/2}\,dx
=
\frac{2!}{(1/2)^3}
=
16.
\]

Equivalently, \(X\) has a gamma distribution with shape \(2\) and scale
\(2\), whose mean is \(2\cdot2=4\).

\item First determine \(c\) by requiring the density to integrate to one:

\begin{align*}
1
&=
\int_{-1}^{1}c(1-x^2)\,dx\\
&=
c\left[x-\frac{x^3}{3}\right]_{-1}^{1}\\
&=
c\left(\frac43\right)\Rightarrow c=\frac{3}{4}
\end{align*}


Now,

\[
E[X]
=
\int_{-1}^{1}x\frac34(1-x^2)\,dx = 0 \ \ \text{(symmetric interval around 0 of an odd function is zero)}.
\]


\item First verify that \(f\) is a density:

\[
\int_{5}^{\infty}\frac{5}{x^2}\,dx
=
5\left[-\frac1x\right]_{5}^{\infty}
=
1.
\]

The expected value is

\begin{align*}
E[X]
&=
\int_{5}^{\infty}x\frac{5}{x^2}\,dx\\
&=
5\int_{5}^{\infty}\frac1x\,dx.
\end{align*}

The integral diverges. Hence, $ {E[X]=\infty}.$



\end{enumerate}

\end{solution}
 


%%%%%%%%%%%%%%%%%%%%%%%%%%%%%%%%%%%%%%%%%%%%%%%%%%%%%%%%%%%%%%
%%%%%%%%%%%%%%%%%%%%%%  PROBLEM 36  %%%%%%%%%%%%%%%%%%%%%%%%%%
%%%%%%%%%%%%%%%%%%%%%%%%%%%%%%%%%%%%%%%%%%%%%%%%%%%%%%%%%%%%%%

   \newpage \begin{problem}

Suppose \(X\) has density

\[
f(x)=
\begin{cases}
a+bx^2, & 0\leq x\leq1,\\
0, & \text{otherwise}.
\end{cases}
\]

If

\[
E[X]=\frac35,
\]

find \(a\) and \(b\).
\end{problem}
\begin{solution}

Because \(f\) is a probability density function,

\[
\int_{0}^{1}(a+bx^2)\,dx=1.
\]

Thus,

\[
a+\frac{b}{3}=1.
\]

Also,

\begin{align*}
E[X]
&=
\int_{0}^{1}x(a+bx^2)\,dx\\
&=
\frac{a}{2}+\frac{b}{4}.
\end{align*}

Since \(E[X]=\frac35\),

\[
\frac{a}{2}+\frac{b}{4}
=
\frac35.
\]

We therefore solve

\[
\begin{cases}
a+\dfrac{b}{3}=1,\\[2mm]
\dfrac{a}{2}+\dfrac{b}{4}=\dfrac35.
\end{cases}
\]

From the first equation,

\[
a=1-\frac{b}{3}.
\]

Substituting into the second equation gives

\begin{align*}
\frac12\left(1-\frac{b}{3}\right)+\frac{b}{4}
&=
\frac35,\\
\frac12+\frac{b}{12}
&=
\frac35,\\
\frac{b}{12}
&=
\frac1{10}.
\end{align*}

Hence,

\[
b=\frac65.
\]

Then

\[
a
=
1-\frac{1}{3}\left(\frac65\right)
=
1-\frac25
=
\frac35.
\]

Therefore,

\[
 {
a=\frac35,
\qquad
b=\frac65.
}
\]

The resulting density is nonnegative on \([0,1]\), as required.

\end{solution}
 


%%%%%%%%%%%%%%%%%%%%%%%%%%%%%%%%%%%%%%%%%%%%%%%%%%%%%%%%%%%%%%
%%%%%%%%%%%%%%%%%%%%%%  PROBLEM 37  %%%%%%%%%%%%%%%%%%%%%%%%%%
%%%%%%%%%%%%%%%%%%%%%%%%%%%%%%%%%%%%%%%%%%%%%%%%%%%%%%%%%%%%%%

   \newpage \begin{problem}

The lifetime \(X\), measured in hours, of an electronic tube has density

\[
f(x)=xe^{-x},
\qquad x\geq0.
\]

Compute the expected lifetime.
\end{problem}
\begin{solution}

The expected lifetime is

\begin{align*}
E[X]
&=
\int_{0}^{\infty}x f(x)\,dx\\
&=
\int_{0}^{\infty}x^2e^{-x}\,dx.
\end{align*}

Using

\[
\int_{0}^{\infty}x^ne^{-x}\,dx=n!,
\]

we obtain

\[
E[X]
=
2!
=
 {2}.
\]

Thus the expected lifetime is

\[
 {2\text{ hours}}.
\]

\end{solution}
 


%%%%%%%%%%%%%%%%%%%%%%%%%%%%%%%%%%%%%%%%%%%%%%%%%%%%%%%%%%%%%%
%%%%%%%%%%%%%%%%%%%%%%  PROBLEM 38  %%%%%%%%%%%%%%%%%%%%%%%%%%
%%%%%%%%%%%%%%%%%%%%%%%%%%%%%%%%%%%%%%%%%%%%%%%%%%%%%%%%%%%%%%

   \newpage \begin{problem} 

\textbf{ *} A point is chosen at random on a line segment of length \(L\).

\begin{enumerate}[label=(\alph*)]
    \item Give a mathematical interpretation of the statement that the
    point is chosen at random.
    \item Find the probability that the ratio of the shorter resulting
    segment to the longer resulting segment is less than \(\frac14\).
\end{enumerate}
\end{problem}
\begin{solution}

Let \(X\) be the distance from the left endpoint of the segment to the
randomly selected point.

Choosing the point uniformly at random means that

\[
X\sim\operatorname{Unif}(0,L),
\]

with density

\[
f_X(x)=\frac1L,
\qquad 0<x<L.
\]

The two resulting segment lengths are

\[
X
\qquad\text{and}\qquad
L-X.
\]

By symmetry, first consider \(0<X\leq L/2\). In this case, \(X\) is the
shorter segment and \(L-X\) is the longer segment. The desired condition is

\[
\frac{X}{L-X}<\frac14.
\]

Solving,

\begin{align*}
4X&<L-X,\\
5X&<L,\\
X&<\frac{L}{5}.
\end{align*}

By symmetry, the same condition occurs near the right endpoint when

\[
X>\frac{4L}{5}.
\]

Thus the favorable set is

\[
\left(0,\frac{L}{5}\right)
\cup
\left(\frac{4L}{5},L\right).
\]

Its total length is

\[
\frac{L}{5}+\frac{L}{5}
=
\frac{2L}{5}.
\]

Therefore,

\[
P\left(
\frac{\text{shorter segment}}{\text{longer segment}}<\frac14
\right)
=
\frac{2L/5}{L}
=
 {\frac25}.
\]

\end{solution}
 


%%%%%%%%%%%%%%%%%%%%%%%%%%%%%%%%%%%%%%%%%%%%%%%%%%%%%%%%%%%%%%
%%%%%%%%%%%%%%%%%%%%%%  PROBLEM 39  %%%%%%%%%%%%%%%%%%%%%%%%%%
%%%%%%%%%%%%%%%%%%%%%%%%%%%%%%%%%%%%%%%%%%%%%%%%%%%%%%%%%%%%%%

   \newpage \begin{problem}

Let

\[
X\sim N(10,36).
\]

Compute the following probabilities:

\begin{enumerate}[label=(\alph*)]
    \item \(P(X>5)\);
    \item \(P(4<X<16)\);
    \item \(P(X<8)\);
    \item \(P(X<20)\);

\end{enumerate}
\end{problem}
\begin{solution}

Since

\[
X\sim N(10,36),
\]

the mean is

\[
\mu=10
\]

and the standard deviation is

\[
\sigma=6.
\]

Define the standardized random variable

\[
Z=\frac{X-10}{6},
\]

so that \(Z\sim N(0,1)\). Let \(\Phi\) denote the standard normal CDF.

\begin{enumerate}[label=(\alph*)]

\item

\begin{align*}
P(X>5)
&=
P\left(
Z>\frac{5-10}{6}
\right)\\
&=
P\left(Z>-\frac56\right)\\
&=
\Phi\left(\frac56\right)\\
&\approx
 {0.7977}.
\end{align*}

\item

\begin{align*}
P(4<X<16)
&=
P\left(
\frac{4-10}{6}<Z<
\frac{16-10}{6}
\right)\\
&=
P(-1<Z<1)\\
&=
\Phi(1)-\Phi(-1)\\
&\approx
 {0.6827}.
\end{align*}

\item

\begin{align*}
P(X<8)
&=
P\left(
Z<\frac{8-10}{6}
\right)\\
&=
P\left(Z<-\frac13\right)\\
&=
\Phi\left(-\frac13\right)\\
&\approx
 {0.3694}.
\end{align*}

\item

\begin{align*}
P(X<20)
&=
P\left(
Z<\frac{20-10}{6}
\right)\\
&=
P\left(Z<\frac53\right)\\
&=
\Phi\left(\frac53\right) \approx  {0.9522}.
\end{align*}

\item

\begin{align*}
P(X>16)
&=
P\left(
Z>\frac{16-10}{6}
\right)\\
&=
P(Z>1)\\
&=
1-\Phi(1)\\
&\approx
 {0.1587}.
\end{align*}

\end{enumerate}

\end{solution}
 


%%%%%%%%%%%%%%%%%%%%%%%%%%%%%%%%%%%%%%%%%%%%%%%%%%%%%%%%%%%%%%
%%%%%%%%%%%%%%%%%%%%%%  PROBLEM 40  %%%%%%%%%%%%%%%%%%%%%%%%%%
%%%%%%%%%%%%%%%%%%%%%%%%%%%%%%%%%%%%%%%%%%%%%%%%%%%%%%%%%%%%%%

   \newpage \begin{problem}

The salaries of physicians in a certain specialty are approximately
normally distributed. Suppose that

\[
P(X<180{,}000)=0.25
\]

and

\[
P(X>320{,}000)=0.25.
\]

Approximately what fraction of physicians earn

\begin{enumerate}[label=(\alph*)]
    \item less than \(\$200{,}000\)?
    \item between \(\$280{,}000\) and \(\$320{,}000\)?
\end{enumerate}
\end{problem}
\begin{solution}

The lower and upper quartiles are

\[
Q_1=180{,}000,
\qquad
Q_3=320{,}000.
\]

Because the normal distribution is symmetric, its mean is the midpoint:

\[
\mu
=
\frac{180{,}000+320{,}000}{2}
=
250{,}000.
\]

For a standard normal random variable,

\[
\Phi^{-1}(0.75)\approx0.67449.
\]

Thus,

\[
\frac{320{,}000-250{,}000}{\sigma}
=
0.67449.
\]

Therefore,

\[
\sigma
=
\frac{70{,}000}{0.67449}
\approx
103{,}782.
\]

\begin{enumerate}[label=(\alph*)]

\item Standardize \(200{,}000\):

\[
z
=
\frac{200{,}000-250{,}000}{103{,}782}
\approx
-0.4818.
\]

Hence,

\[
\begin{aligned}
P(X<200{,}000)
&=
\Phi(-0.4818)\\
&\approx
 {0.3150}.
\end{aligned}
\]

Thus approximately \(31.5\%\) of physicians earn less than
\(\$200{,}000\).

\item Standardize \(280{,}000\):

\[
z
=
\frac{280{,}000-250{,}000}{103{,}782}
\approx
0.2891.
\]

Also, because \(320{,}000\) is the upper quartile,

\[
P(X<320{,}000)=0.75.
\]

Therefore,

\begin{align*}
P(280{,}000<X<320{,}000)
&=
P(X<320{,}000)-P(X<280{,}000)\\
&=
0.75-\Phi(0.2891)\\
&\approx
0.75-0.6137\\
&=
 {0.1363}.
\end{align*}

Thus approximately \(13.6\%\) of physicians earn between
\(\$280{,}000\) and \(\$320{,}000\).

\end{enumerate}

\end{solution}
 

%%%%%%%%%%%%%%%%%%%%%%%%%%%%%%%%%%%%%%%%%%%%%%%%%%%%%%%%%%%%%%
%%%%%%%%%%%%%%%%%%%%%%  PROBLEM 41  %%%%%%%%%%%%%%%%%%%%%%%%%%
%%%%%%%%%%%%%%%%%%%%%%%%%%%%%%%%%%%%%%%%%%%%%%%%%%%%%%%%%%%%%%

   \newpage \begin{problem}

\textbf{*} Let \(Y\) be a continuous random variable with density \(f_Y\), and assume
that \(E[|Y|]<\infty\). Show that

\[
E[Y]
=
\int_{0}^{\infty}P(Y>y)\,dy
-
\int_{0}^{\infty}P(Y<-y)\,dy.
\]
\end{problem}
\begin{solution}

Write $Y=Y^+-Y^-,$ where $Y^+=\max(Y,0),
\text{ and }
Y^-=\max(-Y,0).$ Then $E[Y]=E[Y^+]-E[Y^-].
$
We first consider the positive part. Since

\[
P(Y>y)
=
\int_y^\infty f_Y(x)\,dx,
\]

Tonelli's theorem gives

\begin{align*}
\int_0^\infty P(Y>y)\,dy
&=
\int_0^\infty
\int_y^\infty f_Y(x)\,dx\,dy\\
&=
\int_0^\infty
\int_0^x dy\, f_Y(x)\,dx\\
&=
\int_0^\infty x f_Y(x)\,dx\\
&=
E[Y^+].
\end{align*}

Similarly,

\[
P(Y<-y)
=
\int_{-\infty}^{-y}f_Y(x)\,dx.
\]

Hence,

\begin{align*}
\int_0^\infty P(Y<-y)\,dy
&=
\int_0^\infty
\int_{-\infty}^{-y}f_Y(x)\,dx\,dy.
\end{align*}

For a fixed \(x<0\), the condition \(x<-y\) is equivalent to

\[
0<y<-x.
\]

Therefore, changing the order of integration gives

\begin{align*}
\int_0^\infty P(Y<-y)\,dy
&=
\int_{-\infty}^{0}
\int_0^{-x}dy\, f_Y(x)\,dx\\
&=
\int_{-\infty}^{0}(-x)f_Y(x)\,dx\\
&=
-\int_{-\infty}^{0}x f_Y(x)\,dx\\
&=
E[Y^-].
\end{align*}

Combining the two identities,

\begin{align*}
E[Y]
&=
E[Y^+]-E[Y^-]\\
&=
\int_{0}^{\infty}P(Y>y)\,dy
-
\int_{0}^{\infty}P(Y<-y)\,dy.
\end{align*}

Thus,

\[
 {
E[Y]
=
\int_{0}^{\infty}P(Y>y)\,dy
-
\int_{0}^{\infty}P(Y<-y)\,dy.
}
\]

\end{solution}
 


%%%%%%%%%%%%%%%%%%%%%%%%%%%%%%%%%%%%%%%%%%%%%%%%%%%%%%%%%%%%%%
%%%%%%%%%%%%%%%%%%%%%%  PROBLEM 42  %%%%%%%%%%%%%%%%%%%%%%%%%%
%%%%%%%%%%%%%%%%%%%%%%%%%%%%%%%%%%%%%%%%%%%%%%%%%%%%%%%%%%%%%%

   \newpage \begin{problem}

Let \(Z\sim N(0,1)\). For \(x>0\), show that

\begin{enumerate}[label=(\alph*)]
    \item
    \[
    P(Z>x)=P(Z<-x);
    \]

    \item
    \[
    P(|Z|>x)=2P(Z>x);
    \]

    \item
    \[
    P(|Z|<x)=2P(Z<x)-1.
    \]
\end{enumerate}
\end{problem}
\begin{solution}

The standard normal density is

\[
\phi(z)
=
\frac{1}{\sqrt{2\pi}}e^{-z^2/2}.
\]

Because

\[
\phi(-z)=\phi(z),
\]

the density is symmetric about zero.

\begin{enumerate}[label=(\alph*)]

\item

\begin{align*}
P(Z<-x)
&=
\int_{-\infty}^{-x}\phi(z)\,dz.
\end{align*}

Using the substitution \(u=-z\),

\begin{align*}
P(Z<-x)
&=
\int_x^\infty \phi(-u)\,du\\
&=
\int_x^\infty \phi(u)\,du\\
&=
P(Z>x).
\end{align*}

Therefore,

\[
 {P(Z>x)=P(Z<-x)}.
\]

\item The events \(\{Z>x\}\) and \(\{Z<-x\}\) are disjoint, so

\begin{align*}
P(|Z|>x)
&=
P(Z>x)+P(Z<-x)\\
&=
2P(Z>x).
\end{align*}

Thus,

\[
 {P(|Z|>x)=2P(Z>x)}.
\]

\item By symmetry,

\[
P(Z\leq -x)=P(Z\geq x)=1-\Phi(x).
\]

Therefore,

\begin{align*}
P(|Z|<x)
&=
P(-x<Z<x)\\
&=
\Phi(x)-\Phi(-x)\\
&=
\Phi(x)-\bigl(1-\Phi(x)\bigr)\\
&=
2\Phi(x)-1.
\end{align*}

Since \(\Phi(x)=P(Z<x)\), $P(|Z|<x)=2P(Z<x)-1.$



\end{enumerate}

\end{solution}



%%%%%%%%%%%%%%%%%%%%%%%%%%%%%%%%%%%%%%%%%%%%%%%%%%%%%%%%%%%%%%
%%%%%%%%%%%%%%%%%%%%%%  PROBLEM 43  %%%%%%%%%%%%%%%%%%%%%%%%%%
%%%%%%%%%%%%%%%%%%%%%%%%%%%%%%%%%%%%%%%%%%%%%%%%%%%%%%%%%%%%%%

   \newpage \begin{problem}

 Let

\[
f(x)
=
\frac{1}{\sigma\sqrt{2\pi}}
\exp\left(
-\frac{(x-\mu)^2}{2\sigma^2}
\right)
\]

be the density of a normal random variable with mean \(\mu\) and variance
\(\sigma^2\). Show that

\[
x=\mu-\sigma
\qquad\text{and}\qquad
x=\mu+\sigma
\]

are inflection points of \(f\).
\end{problem}
\begin{solution}

Differentiate \(f\). By the chain rule,

\begin{align*}
f'(x)
&=
f(x)
\left(
-\frac{x-\mu}{\sigma^2}
\right)\\
&=
-\frac{x-\mu}{\sigma^2}f(x).
\end{align*}

Differentiating again,

\begin{align*}
f''(x)
&=
-\frac{1}{\sigma^2}f(x)
-
\frac{x-\mu}{\sigma^2}f'(x)\\
&=
-\frac{1}{\sigma^2}f(x)
+
\frac{(x-\mu)^2}{\sigma^4}f(x)\\
&=
\frac{(x-\mu)^2-\sigma^2}{\sigma^4}f(x).
\end{align*}

Since \(f(x)>0\) for every \(x\),

\[
f''(x)=0
\]

if and only if

\[
(x-\mu)^2-\sigma^2=0.
\]

Thus,

\[
(x-\mu)^2=\sigma^2,
\]

which gives

\[
x-\mu=\pm\sigma.
\]

Therefore,

\[
 {x=\mu-\sigma\quad\text{or}\quad x=\mu+\sigma}.
\]

Moreover,

\[
f''(x)>0
\]

when \(|x-\mu|>\sigma\), while

\[
f''(x)<0
\]

when \(|x-\mu|<\sigma\). Hence the concavity changes at both points, so
they are indeed inflection points.

\end{solution}
 


%%%%%%%%%%%%%%%%%%%%%%%%%%%%%%%%%%%%%%%%%%%%%%%%%%%%%%%%%%%%%%
%%%%%%%%%%%%%%%%%%%%%%  PROBLEM 44  %%%%%%%%%%%%%%%%%%%%%%%%%%
%%%%%%%%%%%%%%%%%%%%%%%%%%%%%%%%%%%%%%%%%%%%%%%%%%%%%%%%%%%%%%

   \newpage \begin{problem}

The median of a continuous random variable with CDF \(F\) is a value
\(m\) satisfying

\[
F(m)=\frac12.
\]

Find the median when \(X\) is

\begin{enumerate}[label=(\alph*)]
    \item uniformly distributed on \((a,b)\);
    \item normally distributed with parameters \(\mu\) and \(\sigma^2\);
    \item exponentially distributed with rate \(\lambda\).
\end{enumerate}
\end{problem}
\begin{solution}

\begin{enumerate}[label=(\alph*)]

\item If \(X\sim\operatorname{Unif}(a,b)\), then for \(a\leq x\leq b\),

\[
F(x)=\frac{x-a}{b-a}.
\]

Set \(F(m)=1/2\):

\begin{align*}
\frac{m-a}{b-a}
&=
\frac12,\\
m-a
&=
\frac{b-a}{2},\\
m
&=
\frac{a+b}{2}.
\end{align*}

Thus,

\[
 {m=\frac{a+b}{2}}.
\]

\item A normal distribution is symmetric about its mean \(\mu\). Hence

\[
P(X<\mu)=\frac12.
\]

Therefore,

\[
 {m=\mu}.
\]

\item If \(X\sim\operatorname{Exp}(\lambda)\), then for \(x\geq0\),

\[
F(x)=1-e^{-\lambda x}.
\]

Set \(F(m)=1/2\):

\begin{align*}
1-e^{-\lambda m}
&=
\frac12,\\
e^{-\lambda m}
&=
\frac12,\\
-\lambda m
&=
-\log 2.
\end{align*}

Therefore,

\[
 {m=\frac{\log 2}{\lambda}}.
\]

\end{enumerate}

\end{solution}
 


%%%%%%%%%%%%%%%%%%%%%%%%%%%%%%%%%%%%%%%%%%%%%%%%%%%%%%%%%%%%%%
%%%%%%%%%%%%%%%%%%%%%%  PROBLEM 45  %%%%%%%%%%%%%%%%%%%%%%%%%%
%%%%%%%%%%%%%%%%%%%%%%%%%%%%%%%%%%%%%%%%%%%%%%%%%%%%%%%%%%%%%%

   \newpage \begin{problem}

 Let \(X\) be a continuous random variable with continuous cumulative
distribution function \(F\). Define

\[
Y=F(X).
\]

Show that \(Y\) is uniformly distributed on \((0,1)\).
\end{problem}
\begin{solution}

Let \(0<y<1\), and let

\[
F^{-1}(y)
=
\inf\{x:F(x)\geq y\}
\]

denote the generalized inverse of \(F\).

Because \(F\) is continuous,

\[
F(F^{-1}(y))=y.
\]

Now,

\begin{align*}
P(Y\leq y)
&=
P(F(X)\leq y)\\
&=
P(X\leq F^{-1}(y))\\
&=
F(F^{-1}(y))\\
&=
y.
\end{align*}

Therefore, the CDF of \(Y\) is

\[
F_Y(y)
=
\begin{cases}
0, & y\leq0,\\
y, & 0<y<1,\\
1, & y\geq1.
\end{cases}
\]

This is the CDF of a uniform random variable on \((0,1)\). Hence,

\[
 {F(X)\sim\operatorname{Unif}(0,1)}.
\]

\end{solution}
 


%%%%%%%%%%%%%%%%%%%%%%%%%%%%%%%%%%%%%%%%%%%%%%%%%%%%%%%%%%%%%%
%%%%%%%%%%%%%%%%%%%%%%  PROBLEM 46  %%%%%%%%%%%%%%%%%%%%%%%%%%
%%%%%%%%%%%%%%%%%%%%%%%%%%%%%%%%%%%%%%%%%%%%%%%% 


%%%%%%%%%%%%%%%%%%%%%%%%%%%%%%%%%%%%%%%%%%%%%%%%%%%%%%%%%%%%%%
%%%%%%%%%%%%%%%%%%%%%%  PROBLEM 47  %%%%%%%%%%%%%%%%%%%%%%%%%%
%%%%%%%%%%%%%%%%%%%%%%%%%%%%%%%%%%%%%%%%%%%%%%%%%%%%%%%%%%%%%%
 


%%%%%%%%%%%%%%%%%%%%%%%%%%%%%%%%%%%%%%%%%%%%%%%%%%%%%%%%%%%%%%
%%%%%%%%%%%%%%%%%%%%%%  PROBLEM 48  %%%%%%%%%%%%%%%%%%%%%%%%%%
%%%%%%%%%%%%%%%%%%%%%%%%%%%%%%%%%%%%%%%%%%%%%%%%%%%%%%%%%%%%%%

   \newpage \begin{problem}

Let

\[
\Phi(x)=P(Z\leq x),
\]

where \(Z\sim N(0,1)\). Determine which of the following statements are
true:

\begin{enumerate}[label=(\alph*)]
    \item
    \[
    \Phi(-x)=\Phi(x);
    \]

    \item
    \[
    \Phi(x)+\Phi(-x)=1;
    \]

    \item
    \[
    \Phi(-x)=\frac{1}{\Phi(x)}.
    \]
\end{enumerate}
\end{problem}
\begin{solution}

By symmetry of the standard normal distribution,

\[
P(Z\leq -x)=P(Z\geq x).
\]

Since the normal distribution is continuous,

\[
P(Z\geq x)=1-P(Z\leq x)=1-\Phi(x).
\]

Therefore,

\[
 {\Phi(-x)=1-\Phi(x)}.
\]

It follows that:

\begin{enumerate}[label=(\alph*)]

\item The statement

\[
\Phi(-x)=\Phi(x)
\]

is generally false. It holds only when \(x=0\).

\item

\[
\Phi(x)+\Phi(-x)
=
\Phi(x)+1-\Phi(x)
=
1.
\]

Thus this statement is

\[
 {\text{true}}.
\]

\item The statement

\[
\Phi(-x)=\frac{1}{\Phi(x)}
\]

is false. In fact,

\[
\Phi(-x)=1-\Phi(x).
\]

\end{enumerate}

Hence, only statement \( {\text{(b)}}\) is true for every real
\(x\).

\end{solution}
 


%%%%%%%%%%%%%%%%%%%%%%%%%%%%%%%%%%%%%%%%%%%%%%%%%%%%%%%%%%%%%%
%%%%%%%%%%%%%%%%%%%%%%  PROBLEM 49  %%%%%%%%%%%%%%%%%%%%%%%%%%
%%%%%%%%%%%%%%%%%%%%%%%%%%%%%%%%%%%%%%%%%%%%%%%%%%%%%%%%%%%%%%

   \newpage \begin{problem}

Let

\[
f(x)=
\begin{cases}
\dfrac13e^x, & x<0,\\[2mm]
\dfrac13, & 0\leq x<1,\\[2mm]
\dfrac13e^{-(x-1)}, & x\geq1.
\end{cases}
\]

\begin{enumerate}[label=(\alph*)]
    \item Show that \(f\) is a probability density function.
    \item If \(X\) has density \(f\), find \(E[X]\).
\end{enumerate}
\end{problem}
\begin{solution}

\begin{enumerate}[label=(\alph*)]

\item Clearly,

\[
f(x)\geq0
\]

for every \(x\). Also,

\begin{align*}
\int_{-\infty}^{\infty}f(x)\,dx
&=
\frac13\int_{-\infty}^{0}e^x\,dx
+
\frac13\int_{0}^{1}1\,dx
+
\frac13\int_{1}^{\infty}e^{-(x-1)}\,dx\\
&=
\frac13(1)
+
\frac13(1)
+
\frac13(1)\\
&=
1.
\end{align*}

Therefore, \(f\) is a probability density function.

\item

\begin{align*}
E[X]
&=
\frac13\int_{-\infty}^{0}xe^x\,dx
+
\frac13\int_{0}^{1}x\,dx
+
\frac13\int_{1}^{\infty}xe^{-(x-1)}\,dx.
\end{align*}

For the first integral,

\[
\int xe^x\,dx=(x-1)e^x,
\]

so

\[
\int_{-\infty}^{0}xe^x\,dx=-1.
\]

For the second integral,

\[
\int_0^1x\,dx=\frac12.
\]

For the third integral, let \(u=x-1\). Then \(x=u+1\), so

\begin{align*}
\int_{1}^{\infty}xe^{-(x-1)}\,dx
&=
\int_{0}^{\infty}(u+1)e^{-u}\,du\\
&=
\int_0^\infty ue^{-u}\,du
+
\int_0^\infty e^{-u}\,du\\
&=
1+1\\
&=
2.
\end{align*}

Therefore,

\begin{align*}
E[X]
&=
\frac13(-1)
+
\frac13\left(\frac12\right)
+
\frac13(2) =  {\frac12}.
\end{align*}

\end{enumerate}

\end{solution}
 


%%%%%%%%%%%%%%%%%%%%%%%%%%%%%%%%%%%%%%%%%%%%%%%%%%%%%%%%%%%%%%
%%%%%%%%%%%%%%%%%%%%%%  PROBLEM 50  %%%%%%%%%%%%%%%%%%%%%%%%%%
%%%%%%%%%%%%%%%%%%%%%%%%%%%%%%%%%%%%%%%%%%%%%%%%%%%%%%%%%%%%%%

   \newpage \begin{problem}
Let

\[
f(x)
=
\frac{\theta^2}{1+\theta}(1+x)e^{-\theta x},
\qquad x>0,
\]

where \(\theta>0\), and let \(f(x)=0\) for \(x\leq0\).

\begin{enumerate}[label=(\alph*)]
    \item Show that \(f\) is a probability density function.
    \item Find \(E[X]\).
    \item Find \(\operatorname{Var}(X)\).
\end{enumerate}
\end{problem}
\begin{solution}

Since \(\theta>0\), $f(x)\geq0.$ We will use the gamma-integral identity

\[
\int_0^\infty x^k e^{-\theta x}\,dx
=
\frac{k!}{\theta^{k+1}},
\qquad k=0,1,2,\ldots
\]

\begin{enumerate}[label=(\alph*)]

\item

\begin{align*}
\int_0^\infty f(x)\,dx
&=
\frac{\theta^2}{1+\theta}
\int_0^\infty(1+x)e^{-\theta x}\,dx\\
&=
\frac{\theta^2}{1+\theta}
\left(
\frac1\theta+\frac1{\theta^2}
\right)\\
&=
\frac{\theta^2}{1+\theta}
\left(
\frac{\theta+1}{\theta^2}
\right)\\
&=
1.
\end{align*}

\item

\begin{align*}
E[X]
&=
\int_0^\infty x f(x)\,dx\\
&=
\frac{\theta^2}{1+\theta}
\int_0^\infty x(1+x)e^{-\theta x}\,dx\\
&=
\frac{\theta^2}{1+\theta}
\left(
\int_0^\infty xe^{-\theta x}\,dx
+
\int_0^\infty x^2e^{-\theta x}\,dx
\right)\\
&=
\frac{\theta^2}{1+\theta}
\left(
\frac1{\theta^2}
+
\frac2{\theta^3}
\right) =  {
\frac{\theta+2}{\theta(\theta+1)}
}.
\end{align*}

\item First compute the second moment:

\begin{align*}
E[X^2]
&=
\int_0^\infty x^2f(x)\,dx\\
&=
\frac{\theta^2}{1+\theta}
\int_0^\infty x^2(1+x)e^{-\theta x}\,dx\\
&=
\frac{\theta^2}{1+\theta}
\left(
\frac2{\theta^3}
+
\frac6{\theta^4}
\right)\\
&=
\frac{2(\theta+3)}
{\theta^2(\theta+1)}.
\end{align*}

Hence,

\begin{align*}
\operatorname{Var}(X)
&=
E[X^2]-E[X]^2\\
&=
\frac{2(\theta+3)}
{\theta^2(\theta+1)}
-
\left(
\frac{\theta+2}
{\theta(\theta+1)}
\right)^2\\
&=
\frac{
2(\theta+3)(\theta+1)-(\theta+2)^2
}
{\theta^2(\theta+1)^2}\\
&=
\frac{
2\theta^2+8\theta+6
-
(\theta^2+4\theta+4)
}
{\theta^2(\theta+1)^2} =  {
\frac{\theta^2+4\theta+2}
{\theta^2(\theta+1)^2}
}.
\end{align*}

\end{enumerate}

\end{solution}
 











%%%%%%%%%%%%%%%%%%%%%%%%%%%%%%%%%%%%%%%%%%%%%%%%%%%%%%%%%%%%%%
%%%%%%%%%%%%%%%%%%%%%%  PROBLEM 51  %%%%%%%%%%%%%%%%%%%%%%%%%%
%%%%%%%%%%%%%%%%%%%%%%%%%%%%%%%%%%%%%%%%%%%%%%%%%%%%%%%%%%%%%%
 


%%%%%%%%%%%%%%%%%%%%%%%%%%%%%%%%%%%%%%%%%%%%%%%%%%%%%%%%%%%%%%
%%%%%%%%%%%%%%%%%%%%%%  PROBLEM 52  %%%%%%%%%%%%%%%%%%%%%%%%%%
%%%%%%%%%%%%%%%%%%%%%%%%%%%%%%%%%%%%%%%%%%%%%%%%%%%%%%%%%%%%%%
 


%%%%%%%%%%%%%%%%%%%%%%%%%%%%%%%%%%%%%%%%%%%%%%%%%%%%%%%%%%%%%%
%%%%%%%%%%%%%%%%%%%%%%  PROBLEM 53  %%%%%%%%%%%%%%%%%%%%%%%%%%
%%%%%%%%%%%%%%%%%%%%%%%%%%%%%%%%%%%%%%%%%%%%%%%%%%%%%%%%%%%%%%

   \newpage \begin{problem}

\textbf{* } Let \(A_1,A_2,\ldots\) be a sequence of events. Prove that

\[
P\left(\bigcup_{n=1}^{\infty}A_n\right)
\leq
\sum_{n=1}^{\infty}P(A_n).
\]
\end{problem}
\begin{solution}

Define disjoint events

\[
B_1=A_1
\]

and, for \(n\geq2\),

\[
B_n
=
A_n\setminus\bigcup_{j=1}^{n-1}A_j.
\]

Then the events \(B_1,B_2,\ldots\) are pairwise disjoint and

\[
\bigcup_{n=1}^{\infty}B_n
=
\bigcup_{n=1}^{\infty}A_n.
\]

Also,

\[
B_n\subseteq A_n,
\]

so

\[
P(B_n)\leq P(A_n).
\]

By countable additivity,

\begin{align*}
P\left(\bigcup_{n=1}^{\infty}A_n\right)
&=
P\left(\bigcup_{n=1}^{\infty}B_n\right)\\
&=
\sum_{n=1}^{\infty}P(B_n)\\
&\leq
\sum_{n=1}^{\infty}P(A_n).
\end{align*}

Thus,

\[
 {
P\left(\bigcup_{n=1}^{\infty}A_n\right)
\leq
\sum_{n=1}^{\infty}P(A_n).
}
\]

\end{solution}


%%%%%%%%%%%%%%%%%%%%%%%%%%%%%%%%%%%%%%%%%%%%%%%%%%%%%%%%%%%%%%
%%%%%%%%%%%%%%%%%%%%%%  PROBLEM 54  %%%%%%%%%%%%%%%%%%%%%%%%%%
%%%%%%%%%%%%%%%%%%%%%%%%%%%%%%%%%%%%%%%%%%%%%%%%%%%%%%%%%%%%%%

   \newpage \begin{problem}

\textbf{* }Let \(A_1,A_2,\ldots\) be a sequence of events. Prove that

\[
P\left(\bigcap_{n=1}^{\infty}A_n\right)
\geq
1-\sum_{n=1}^{\infty}P(A_n^c).
\]
\end{problem}
\begin{solution}

By De Morgan's law,

\[
\left(\bigcap_{n=1}^{\infty}A_n\right)^c
=
\bigcup_{n=1}^{\infty}A_n^c.
\]

Therefore,

\begin{align*}
P\left(\bigcap_{n=1}^{\infty}A_n\right)
&=
1-
P\left(
\bigcup_{n=1}^{\infty}A_n^c
\right).
\end{align*}

By Boole's inequality,

\[
P\left(
\bigcup_{n=1}^{\infty}A_n^c
\right)
\leq
\sum_{n=1}^{\infty}P(A_n^c).
\]

Hence,

\[
 {
P\left(\bigcap_{n=1}^{\infty}A_n\right)
\geq
1-\sum_{n=1}^{\infty}P(A_n^c).
}
\]

\end{solution}
 


%%%%%%%%%%%%%%%%%%%%%%%%%%%%%%%%%%%%%%%%%%%%%%%%%%%%%%%%%%%%%%
%%%%%%%%%%%%%%%%%%%%%%  PROBLEM 55  %%%%%%%%%%%%%%%%%%%%%%%%%%
%%%%%%%%%%%%%%%%%%%%%%%%%%%%%%%%%%%%%%%%%%%%%%%%%%%%%%%%%%%%%%

   \newpage \begin{problem}

Determine whether each statement is true. If it is true, prove it. If it
is false, give a counterexample.

\begin{enumerate}[label=(\alph*)]
    \item If \(P(A)=1\), then \(A=\Omega\).
    \item If \(P(B)=0\), then \(B=\varnothing\).
\end{enumerate}
\end{problem}
\begin{solution}

Both statements are false in general.

\begin{enumerate}[label=(\alph*)]

\item Consider

\[
\Omega=[0,1]
\]

with the uniform probability distribution, and let

\[
A=(0,1].
\]

Then

\[
A\neq\Omega,
\]

because \(0\notin A\), but

\[
P(A)=1.
\]

Thus, an event may have probability \(1\) without being the entire sample
space.

\item In the same probability space, let

\[
B=\left\{\frac12\right\}.
\]

Then

\[
B\neq\varnothing,
\]

but a singleton has probability zero under the uniform distribution:

\[
P(B)=0.
\]

Thus, a nonempty event may have probability zero.

\end{enumerate}

The statements would be true in certain finite probability spaces in
which every individual outcome has positive probability, but they are
not true for arbitrary probability spaces.

\end{solution}
 


%%%%%%%%%%%%%%%%%%%%%%%%%%%%%%%%%%%%%%%%%%%%%%%%%%%%%%%%%%%%%%
%%%%%%%%%%%%%%%%%%%%%%  PROBLEM 56  %%%%%%%%%%%%%%%%%%%%%%%%%%
%%%%%%%%%%%%%%%%%%%%%%%%%%%%%%%%%%%%%%%%%%%%%%%%%%%%%%%%%%%%%%

   \newpage \begin{problem}

Let \(A\) and \(B\) be events satisfying

\[
P(A)=P(B)=1.
\]

Show that

\[
P(A\cap B)=1.
\]
\end{problem}
\begin{solution}

By Bonferroni's inequality,

\[
P(A\cap B)
\geq
P(A)+P(B)-1.
\]

Therefore,

\[
P(A\cap B)
\geq
1+1-1
=
1.
\]

Since no probability can exceed \(1\),

\[
 {P(A\cap B)=1}.
\]

Alternatively,

\[
(A\cap B)^c=A^c\cup B^c.
\]

Since

\[
P(A^c)=P(B^c)=0,
\]

Boole's inequality gives

\[
P((A\cap B)^c)
\leq
P(A^c)+P(B^c)
=
0.
\]

Hence \(P(A\cap B)=1\).

\end{solution}
 


%%%%%%%%%%%%%%%%%%%%%%%%%%%%%%%%%%%%%%%%%%%%%%%%%%%%%%%%%%%%%%
%%%%%%%%%%%%%%%%%%%%%%  PROBLEM 57  %%%%%%%%%%%%%%%%%%%%%%%%%%
%%%%%%%%%%%%%%%%%%%%%%%%%%%%%%%%%%%%%%%%%%%%%%%%%%%%%%%%%%%%%%

   \newpage \begin{problem}

Is it possible to define a probability measure on a countably infinite
sample space so that all outcomes are equally probable?
\end{problem}
\begin{solution}

Let

\[
\Omega=\{\omega_1,\omega_2,\ldots\}
\]

be countably infinite, and suppose every outcome has the same probability
\(p\).

If \(p>0\), then by countable additivity,

\[
P(\Omega)
=
\sum_{n=1}^{\infty}P(\{\omega_n\})
=
\sum_{n=1}^{\infty}p
=
\infty,
\]

which contradicts \(P(\Omega)=1\).

If \(p=0\), then

\[
P(\Omega)
=
\sum_{n=1}^{\infty}0
=
0,
\]

which again contradicts \(P(\Omega)=1\).

Therefore,

No countably infinite sample space can have equally probable
outcomes under a countably additive probability measure

\end{solution}
 


%%%%%%%%%%%%%%%%%%%%%%%%%%%%%%%%%%%%%%%%%%%%%%%%%%%%%%%%%%%%%%
%%%%%%%%%%%%%%%%%%%%%%  PROBLEM 58  %%%%%%%%%%%%%%%%%%%%%%%%%%
%%%%%%%%%%%%%%%%%%%%%%%%%%%%%%%%%%%%%%%%%%%%%%%%%%%%%%%%%%%%%%

   \newpage \begin{problem}
Let \(A_1,\ldots,A_n\) be events satisfying

\[
P(A_1)=\cdots=P(A_n)=1.
\]

Show that

\[
P(A_1\cap\cdots\cap A_n)=1.
\]
\end{problem}
\begin{solution}

By De Morgan's law,

\[
(A_1\cap\cdots\cap A_n)^c
=
A_1^c\cup\cdots\cup A_n^c.
\]

Since \(P(A_i)=1\),

\[
P(A_i^c)=0
\]

for every \(i\).

By Boole's inequality,

\[
P(A_1^c\cup\cdots\cup A_n^c)
\leq
\sum_{i=1}^{n}P(A_i^c)
=
0.
\]

Thus,

\[
P((A_1\cap\cdots\cap A_n)^c)=0,
\]

and consequently,

\[
 {
P(A_1\cap\cdots\cap A_n)=1.
}
\]

\end{solution}
 


%%%%%%%%%%%%%%%%%%%%%%%%%%%%%%%%%%%%%%%%%%%%%%%%%%%%%%%%%%%%%%
%%%%%%%%%%%%%%%%%%%%%%  PROBLEM 59  %%%%%%%%%%%%%%%%%%%%%%%%%%
%%%%%%%%%%%%%%%%%%%%%%%%%%%%%%%%%%%%%%%%%%%%%%%%%%%%%%%%%%%%%%

   \newpage \begin{problem}

How many \(n\times m\) matrices with entries in \(\{0,1\}\) are there?
\end{problem}
\begin{solution}

An \(n\times m\) matrix contains

\[
nm
\]

entries. Each entry can independently be chosen to be either \(0\) or
\(1\), giving two choices per entry.

By the multiplication rule, the number of matrices is

\[
 {2^{nm}}.
\]

\end{solution}
 


%%%%%%%%%%%%%%%%%%%%%%%%%%%%%%%%%%%%%%%%%%%%%%%%%%%%%%%%%%%%%%
%%%%%%%%%%%%%%%%%%%%%%  PROBLEM 60  %%%%%%%%%%%%%%%%%%%%%%%%%%
%%%%%%%%%%%%%%%%%%%%%%%%%%%%%%%%%%%%%%%%%%%%%%%%%%%%%%%%%%%%%%

   \newpage \begin{problem}

How many four-digit numbers can be formed using only the digits

\[
2,4,6,8,9?
\]

How many of these numbers have at least one repeated digit?
\end{problem}
\begin{solution}

There are five choices for each of the four digit positions. Since none of
the available digits is zero, every resulting string is a four-digit
number.

Thus, the total number is

\[
5^4=625.
\]

To count numbers with at least one repeated digit, first count those with
all four digits distinct.

There are

\[
5\cdot4\cdot3\cdot2
=
120
\]

such numbers.

Therefore, the number with at least one repeated digit is

\[
625-120
=
 {505}.
\]

Hence,

\[
 {625\text{ total numbers},\qquad505\text{ with repetition}.}
\]

\end{solution}
 


%%%%%%%%%%%%%%%%%%%%%%%%%%%%%%%%%%%%%%%%%%%%%%%%%%%%%%%%%%%%%%
%%%%%%%%%%%%%%%%%%%%%%  PROBLEM 61  %%%%%%%%%%%%%%%%%%%%%%%%%%
%%%%%%%%%%%%%%%%%%%%%%%%%%%%%%%%%%%%%%%%%%%%%%%%%%%%%%%%%%%%%%

   \newpage \begin{problem}

Six fair dice are tossed. Find the probability that at least two of them
show the same face.
\end{problem}
\begin{solution}

Let \(A\) be the event that at least two dice show the same face.

Its complement \(A^c\) is the event that all six dice show different
faces.

There are

\[
6^6
\]

equally likely ordered outcomes.

For all six faces to be different, every face \(1,\ldots,6\) must appear
exactly once. There are

\[
6!
\]

such outcomes.

Therefore,

\[
\begin{aligned}
P(A)
&=
1-P(A^c)\\
&=
1-\frac{6!}{6^6}\\
&=
1-\frac{720}{46656}\\
&=
 {\frac{899}{900}}.
\end{aligned}
\]

\end{solution}
 


%%%%%%%%%%%%%%%%%%%%%%%%%%%%%%%%%%%%%%%%%%%%%%%%%%%%%%%%%%%%%%
%%%%%%%%%%%%%%%%%%%%%%  PROBLEM 62  %%%%%%%%%%%%%%%%%%%%%%%%%%
%%%%%%%%%%%%%%%%%%%%%%%%%%%%%%%%%%%%%%%%%%%%%%%%%%%%%%%%%%%%%%

   \newpage \begin{problem}
Show that if

\[
P(A)=1,
\]

then

\[
P(B\mid A)=P(B).
\]
\end{problem}
\begin{solution}

Since \(P(A)=1\),

\[
P(A^c)=0.
\]

Now,

\[
B=(B\cap A)\cup(B\cap A^c),
\]

where the union is disjoint. Therefore,

\[
P(B)
=
P(B\cap A)+P(B\cap A^c).
\]

Since

\[
B\cap A^c\subseteq A^c,
\]

we have

\[
P(B\cap A^c)=0.
\]

Thus,

\[
P(B\cap A)=P(B).
\]

Because \(P(A)=1>0\),

\[
P(B\mid A)
=
\frac{P(B\cap A)}{P(A)}
=
\frac{P(B)}{1}
=
 {P(B)}.
\]

\end{solution}
 


%%%%%%%%%%%%%%%%%%%%%%%%%%%%%%%%%%%%%%%%%%%%%%%%%%%%%%%%%%%%%%
%%%%%%%%%%%%%%%%%%%%%%  PROBLEM 63  %%%%%%%%%%%%%%%%%%%%%%%%%%
%%%%%%%%%%%%%%%%%%%%%%%%%%%%%%%%%%%%%%%%%%%%%%%%%%%%%%%%%%%%%%

   \newpage \begin{problem}

\textbf{* }An integer is selected uniformly at random from

\[
\{1,2,\ldots,10{,}000\}
\]

and is observed to be odd. Find the conditional probability that it is

\begin{enumerate}[label=(\alph*)]
    \item divisible by \(3\);
    \item divisible by neither \(3\) nor \(5\).
\end{enumerate}
\end{problem}
\begin{solution}

There are exactly

\[
5000
\]

odd integers between \(1\) and \(10{,}000\).

\begin{enumerate}[label=(\alph*)]

\item An odd integer divisible by \(3\) is an odd multiple of \(3\).

There are

\[
\left\lfloor\frac{10{,}000}{3}\right\rfloor
=
3333
\]

multiples of \(3\). Among these, the odd multiples correspond to odd
multipliers \(1,3,\ldots,3333\), of which there are

\[
1667.
\]

Therefore,

\[
 {
P(3\mid\text{odd})
=
\frac{1667}{5000}.
}
\]

\item Among the \(5000\) odd integers, count those divisible by \(3\) or
\(5\).

The number of odd multiples of \(3\) is \(1667\).

There are

\[
\left\lfloor\frac{10{,}000}{5}\right\rfloor=2000
\]

multiples of \(5\), of which \(1000\) are odd.

The numbers divisible by both \(3\) and \(5\) are multiples of \(15\).
There are

\[
\left\lfloor\frac{10{,}000}{15}\right\rfloor=666
\]

multiples of \(15\), of which \(333\) are odd.

By inclusion--exclusion, the number of odd integers divisible by \(3\) or
\(5\) is

\[
1667+1000-333=2334.
\]

Thus, the number divisible by neither is

\[
5000-2334=2666.
\]

Therefore,

\[
 {
P(\text{neither }3\text{ nor }5\mid\text{odd})
=
\frac{2666}{5000}
=
\frac{1333}{2500}.
}
\]

\end{enumerate}

\end{solution}
 


%%%%%%%%%%%%%%%%%%%%%%%%%%%%%%%%%%%%%%%%%%%%%%%%%%%%%%%%%%%%%%
%%%%%%%%%%%%%%%%%%%%%%  PROBLEM 64  %%%%%%%%%%%%%%%%%%%%%%%%%%
%%%%%%%%%%%%%%%%%%%%%%%%%%%%%%%%%%%%%%%%%%%%%%%%%%%%%%%%%%%%%%

   \newpage \begin{problem}

A judge initially assigns probability \(0.65\) to the event that Susan is
guilty. Robert knows whether Susan is guilty or innocent. If Susan is
guilty, Robert lies with probability \(0.25\). If Susan is innocent,
Robert tells the truth. Find the probability that Robert commits perjury.
\end{problem}
\begin{solution}

Let \(G\) be the event that Susan is guilty and \(L\) the event that
Robert lies.

We are given

\[
P(G)=0.65,
\]

\[
P(L\mid G)=0.25,
\]

and

\[
P(L\mid G^c)=0.
\]

By the law of total probability,

\begin{align*}
P(L)
&=
P(L\mid G)P(G)
+
P(L\mid G^c)P(G^c)\\
&=
(0.25)(0.65)
+
(0)(0.35)\\
&=
 {0.1625}.
\end{align*}

Thus, the probability that Robert commits perjury is

\[
 {16.25\%}.
\]

\end{solution}
 


%%%%%%%%%%%%%%%%%%%%%%%%%%%%%%%%%%%%%%%%%%%%%%%%%%%%%%%%%%%%%%
%%%%%%%%%%%%%%%%%%%%%%  PROBLEM 65  %%%%%%%%%%%%%%%%%%%%%%%%%%
%%%%%%%%%%%%%%%%%%%%%%%%%%%%%%%%%%%%%%%%%%%%%%%%%%%%%%%%%%%%%%

   \newpage \begin{problem}

\textbf{* }Let \(X\) denote the time until a new car breaks down, and define

\[
Y=
\begin{cases}
X, & X\leq5,\\
5, & X>5.
\end{cases}
\]

Let \(F\) be the cumulative distribution function of \(X\). Find the
cumulative distribution function of \(Y\) in terms of \(F\).
\end{problem}
\begin{solution}

The random variable \(Y\) can be written as

\[
Y=\min(X,5).
\]

Let

\[
F_Y(y)=P(Y\leq y).
\]

If \(y<5\), then

\[
\{Y\leq y\}=\{X\leq y\},
\]

because whenever \(X>5\), we have \(Y=5>y\). Therefore,

\[
F_Y(y)=F(y),
\qquad y<5.
\]

If \(y\geq5\), then \(Y\leq5\leq y\) for every outcome, so

\[
F_Y(y)=1.
\]

Hence,

\[
 {
F_Y(y)=
\begin{cases}
F(y), & y<5,\\[1mm]
1, & y\geq5.
\end{cases}
}
\]

Notice that \(Y\) has a point mass at \(5\) of size

\[
P(Y=5)
=
1-F(5^-).
\]

This mass includes both the event \(X=5\) and the event \(X>5\).

\end{solution}
 


%%%%%%%%%%%%%%%%%%%%%%%%%%%%%%%%%%%%%%%%%%%%%%%%%%%%%%%%%%%%%%
%%%%%%%%%%%%%%%%%%%%%%  PROBLEM 66  %%%%%%%%%%%%%%%%%%%%%%%%%%
%%%%%%%%%%%%%%%%%%%%%%%%%%%%%%%%%%%%%%%%%%%%%%%%%%%%%%%%%%%%%%

   \newpage \begin{problem}

Suppose the cumulative distribution function of \(X\) is

\[
F(x)=
\begin{cases}
0, & x<-2,\\[1mm]
\dfrac12, & -2\leq x<2,\\[1mm]
\dfrac35, & 2\leq x<4,\\[1mm]
\dfrac89, & 4\leq x<6,\\[1mm]
1, & x\geq6.
\end{cases}
\]

Determine the probability mass function of \(X\) and describe its graph.
\end{problem}
\begin{solution}

For a discrete random variable, the probability at a point is the size
of the jump of the CDF:

\[
P(X=x)=F(x)-F(x^-).
\]

At \(x=-2\),

\[
P(X=-2)
=
\frac12-0
=
\frac12.
\]

At \(x=2\),

\[
P(X=2)
=
\frac35-\frac12
=
\frac1{10}.
\]

At \(x=4\),

\[
P(X=4)
=
\frac89-\frac35
=
\frac{40-27}{45}
=
\frac{13}{45}.
\]

At \(x=6\),

\[
P(X=6)
=
1-\frac89
=
\frac19.
\]

Therefore,

\[
 {
p_X(x)=
\begin{cases}
\dfrac12, & x=-2,\\[1mm]
\dfrac1{10}, & x=2,\\[1mm]
\dfrac{13}{45}, & x=4,\\[1mm]
\dfrac19, & x=6,\\[1mm]
0, & \text{otherwise}.
\end{cases}
}
\]

The probabilities sum to

\[
\frac12+\frac1{10}+\frac{13}{45}+\frac19
=
1.
\]


\end{solution}



%%%%%%%%%%%%%%%%%%%%%%%%%%%%%%%%%%%%%%%%%%%%%%%%%%%%%%%%%%%%%%
%%%%%%%%%%%%%%%%%%%%%%  PROBLEM 67  %%%%%%%%%%%%%%%%%%%%%%%%%%
%%%%%%%%%%%%%%%%%%%%%%%%%%%%%%%%%%%%%%%%%%%%%%%%%%%%%%%%%%%%%%

   \newpage \begin{problem}

For each of the following functions, determine the value of \(k\) for
which \(p\) is a probability mass function.

\begin{enumerate}[label=(\alph*)]

\item
\[
p(x)=kx,
\qquad x=1,2,3,4,5.
\]

\item
\[
p(x)=k(1+x)^2,
\qquad x=-2,0,1,2.
\]

\item
\[
p(x)=k\left(\frac19\right)^x,
\qquad x=1,2,3,\ldots
\]

\item
\[
p(x)=kx,
\qquad x=1,2,\ldots,n.
\]

\item
\[
p(x)=kx^2,
\qquad x=1,2,\ldots,n.
\]

\end{enumerate}
\end{problem}
\begin{solution}

\begin{enumerate}[label=(\alph*)]

\item

\[
1
=
k\sum_{x=1}^{5}x
=
k(1+2+3+4+5)
=
15k\Rightarrow  {k=\frac1{15}}.
\]

 

\item The values of \((1+x)^2\) are

\[
1,\quad1,\quad4,\quad9
\]

for \(x=-2,0,1,2\), respectively. Hence,

\[
1
=
k(1+1+4+9)
=
15k\Rightarrow  {k=\frac1{15}}.
\]

 

\item Using the geometric-series formula,

\begin{align*}
1
&=
k\sum_{x=1}^{\infty}\left(\frac19\right)^x\\
&=
k\frac{1/9}{1-1/9}\\
&=
\frac{k}{8} \Rightarrow k=8
\end{align*}


\item Since

\[
\sum_{x=1}^{n}x
=
\frac{n(n+1)}{2},
\]

we require

\[
1
=
k\frac{n(n+1)}{2} \Rightarrow k=\frac{2}{n(n+1)}.
\]


\item Since

\[
\sum_{x=1}^{n}x^2
=
\frac{n(n+1)(2n+1)}{6},
\]

we require

\[
1
=
k\frac{n(n+1)(2n+1)}{6}.
\]

Hence,

\[
 {
k=\frac{6}{n(n+1)(2n+1)}.
}
\]

\end{enumerate}

\end{solution}
 


%%%%%%%%%%%%%%%%%%%%%%%%%%%%%%%%%%%%%%%%%%%%%%%%%%%%%%%%%%%%%%
%%%%%%%%%%%%%%%%%%%%%%  PROBLEM 68  %%%%%%%%%%%%%%%%%%%%%%%%%%
%%%%%%%%%%%%%%%%%%%%%%%%%%%%%%%%%%%%%%%%%%%%%%%%%%%%%%%%%%%%%%

   \newpage \begin{problem}

Let \(X\) have probability mass function

\[
p(x)=
\begin{cases}
\dfrac{|x-3|+1}{28},
&x=-3,-2,-1,0,1,2,3,\\[2mm]
0,
&\text{otherwise}.
\end{cases}
\]

Find \(\operatorname{Var}(X)\).
\end{problem}
\begin{solution}

The probabilities are

\[
\begin{array}{c|rrrrrrr}
x&-3&-2&-1&0&1&2&3\\ \hline
28p(x)&7&6&5&4&3&2&1.
\end{array}
\]

First compute the mean:

\begin{align*}
E[X]
&=
\frac1{28}
\left[
(-3)(7)+(-2)(6)+(-1)(5)
+0(4)+1(3)+2(2)+3(1)
\right]\\
&=
\frac{-28}{28}\\
&=
-1.
\end{align*}

Next,

\begin{align*}
E[X^2]
&=
\frac1{28}
\left[
9(7)+4(6)+1(5)+0(4)+1(3)+4(2)+9(1)
\right]\\
&=
\frac{112}{28}\\
&=
4.
\end{align*}

Therefore,

\begin{align*}
\operatorname{Var}(X)
&=
E[X^2]-E[X]^2\\
&=
4-(-1)^2\\
&=
 {3}.
\end{align*}

\end{solution}
 


%%%%%%%%%%%%%%%%%%%%%%%%%%%%%%%%%%%%%%%%%%%%%%%%%%%%%%%%%%%%%%
%%%%%%%%%%%%%%%%%%%%%%  PROBLEM 69  %%%%%%%%%%%%%%%%%%%%%%%%%%
%%%%%%%%%%%%%%%%%%%%%%%%%%%%%%%%%%%%%%%%%%%%%%%%%%%%%%%%%%%%%%

   \newpage \begin{problem}

Let \(X\) have cumulative distribution function

\[
F(x)=
\begin{cases}
0, & x<-3,\\[1mm]
\dfrac38, & -3\leq x<0,\\[1mm]
\dfrac34, & 0\leq x<6,\\[1mm]
1, & x\geq6.
\end{cases}
\]

Find the variance and standard deviation of \(X\).
\end{problem}
\begin{solution}

The probability masses are the jumps of \(F\):

\[
P(X=-3)=\frac38,
\]

\[
P(X=0)=\frac34-\frac38=\frac38,
\]

and

\[
P(X=6)=1-\frac34=\frac14.
\]

Thus,

\begin{align*}
E[X]
&=
(-3)\frac38
+
0\frac38
+
6\frac14\\
&=
-\frac98+\frac{12}{8}\\
&=
\frac38.
\end{align*}

Also,

\begin{align*}
E[X^2]
&=
9\frac38
+
0
+
36\frac14\\
&=
\frac{27}{8}+9\\
&=
\frac{99}{8}.
\end{align*}

Therefore,

\begin{align*}
\operatorname{Var}(X)
&=
E[X^2]-E[X]^2\\
&=
\frac{99}{8}
-
\left(\frac38\right)^2\\
&=
\frac{792}{64}-\frac9{64}\\
&=
 {\frac{783}{64}}.
\end{align*}

The standard deviation is

\[
 {
\sigma_X
=
\sqrt{\frac{783}{64}}
=
\frac{\sqrt{783}}{8}
\approx3.4978.
}
\]

\end{solution}
 


%%%%%%%%%%%%%%%%%%%%%%%%%%%%%%%%%%%%%%%%%%%%%%%%%%%%%%%%%%%%%%
%%%%%%%%%%%%%%%%%%%%%%  PROBLEM 70  %%%%%%%%%%%%%%%%%%%%%%%%%%
%%%%%%%%%%%%%%%%%%%%%%%%%%%%%%%%%%%%%%%%%%%%%%%%%%%%%%%%%%%%%%

   \newpage \begin{problem}

Suppose \(X\) is a discrete random variable satisfying

\[
E[X]=1
\]

and

\[
E[X(X-2)]=3.
\]

Find

\[
\operatorname{Var}(-3X+5).
\]
\end{problem}
\begin{solution}

Since

\[
X(X-2)=X^2-2X,
\]

we have

\[
E[X^2]-2E[X]=3.
\]

Using \(E[X]=1\),

\[
E[X^2]-2=3,
\]

so

\[
E[X^2]=5.
\]

Therefore,

\[
\operatorname{Var}(X)
=
E[X^2]-E[X]^2
=
5-1
=
4.
\]

Now use

\[
\operatorname{Var}(aX+b)
=
a^2\operatorname{Var}(X).
\]

Thus,

\[
\operatorname{Var}(-3X+5)
=
(-3)^2(4)
=
 {36}.
\]

\end{solution}
 


%%%%%%%%%%%%%%%%%%%%%%%%%%%%%%%%%%%%%%%%%%%%%%%%%%%%%%%%%%%%%%
%%%%%%%%%%%%%%%%%%%%%%  PROBLEM 71  %%%%%%%%%%%%%%%%%%%%%%%%%%
%%%%%%%%%%%%%%%%%%%%%%%%%%%%%%%%%%%%%%%%%%%%%%%%%%%%%%%%%%%%%%

   \newpage \begin{problem}

Let \(X\) be the amount, in fluid ounces, in a randomly selected bottle
from company \(A\), and let \(Y\) be the amount in a randomly selected
bottle from company \(B\). Their distributions are

\[
\begin{array}{c|ccccc}
x&15.85&15.9&16&16.1&16.2\\ \hline
P(X=x)&0.15&0.21&0.35&0.15&0.14\\
P(Y=x)&0.14&0.05&0.64&0.08&0.09
\end{array}
\]

Find \(E[X]\), \(E[Y]\), \(\operatorname{Var}(X)\), and
\(\operatorname{Var}(Y)\), and interpret the results.
\end{problem}
\begin{solution}

For \(X\),

\begin{align*}
E[X]
&=
15.85(0.15)+15.9(0.21)+16(0.35)\\
&\qquad
+16.1(0.15)+16.2(0.14)\\
&=
 {15.9995}.
\end{align*}

Similarly,

\begin{align*}
E[Y]
&=
15.85(0.14)+15.9(0.05)+16(0.64)\\
&\qquad
+16.1(0.08)+16.2(0.09)\\
&=
 {16.0000}.
\end{align*}

The second moments are

\[
E[X^2]
=
\sum_xx^2P(X=x)
=
256. - \text{computed through the distribution},
\]

and

\[
E[Y^2]
=
\sum_xx^2P(Y=x).
\]

Using

\[
\operatorname{Var}(X)=E[X^2]-E[X]^2,
\]

the resulting variances are

\[
 {
\operatorname{Var}(X)=0.01257475
}
\]

and

\[
 {
\operatorname{Var}(Y)=0.00805.
}
\]

Thus, the standard deviations are approximately

\[
\sigma_X\approx0.1121
\]

and

\[
\sigma_Y\approx0.0897.
\]

Both companies fill their bottles with an average of approximately
\(16\) fluid ounces. Company \(B\), however, has the smaller variance and
standard deviation, so its bottle amounts are more tightly concentrated
around \(16\) ounces. Company \(B\) is therefore more consistent.

\end{solution}
 


%%%%%%%%%%%%%%%%%%%%%%%%%%%%%%%%%%%%%%%%%%%%%%%%%%%%%%%%%%%%%%
%%%%%%%%%%%%%%%%%%%%%%  PROBLEM 72  %%%%%%%%%%%%%%%%%%%%%%%%%%
%%%%%%%%%%%%%%%%%%%%%%%%%%%%%%%%%%%%%%%%%%%%%%%%%%%%%%%%%%%%%%

   \newpage \begin{problem}

A couple wants to have at least a \(95\%\) probability of having at least
one boy and at least one girl. Assume that the sexes of the children are
independent and that a child is equally likely to be a boy or a girl.

What is the minimum number of children they should plan to have?
\end{problem}
\begin{solution}

For \(n\) children, the complement of having at least one boy and at
least one girl is the event that all children have the same sex.

Thus,

\begin{align*}
P(\text{at least one boy and one girl})
&=
1-P(\text{all boys or all girls})\\
&=
1-\left(\frac12\right)^n-\left(\frac12\right)^n\\
&=
1-2^{1-n}.
\end{align*}

We require

\[
1-2^{1-n}\geq0.95.
\]

Therefore,

\[
2^{1-n}\leq0.05.
\]

Testing integers,

\[
n=5:
\qquad
1-2^{1-5}
=
1-\frac1{16}
=
0.9375<0.95,
\]

whereas

\[
n=6:
\qquad
1-2^{1-6}
=
1-\frac1{32}
=
0.96875>0.95.
\]

Hence, the minimum number is

\[
 {6}.
\]

\end{solution}
 


%%%%%%%%%%%%%%%%%%%%%%%%%%%%%%%%%%%%%%%%%%%%%%%%%%%%%%%%%%%%%%
%%%%%%%%%%%%%%%%%%%%%%  PROBLEM 73  %%%%%%%%%%%%%%%%%%%%%%%%%%
%%%%%%%%%%%%%%%%%%%%%%%%%%%%%%%%%%%%%%%%%%%%%%%%%%%%%%%%%%%%%%

   \newpage \begin{problem}

A rare blood type occurs in \(0.05\%\) of the population. A group of
\(3000\) people is selected independently from the population. Find the
probability that at least two people in the group have the rare blood
type.
\end{problem}
\begin{solution}

Let \(X\) be the number of people with the rare blood type. Then

\[
X\sim\operatorname{Bin}(3000,0.0005).
\]

The desired probability is

\begin{align*}
P(X\geq2)
&=
1-P(X=0)-P(X=1)\\
&=
1-(1-0.0005)^{3000}\\
&\qquad
-3000(0.0005)(1-0.0005)^{2999}.
\end{align*}

Thus,

\[
 {
P(X\geq2)
\approx0.4422.
}
\]

Because \(n\) is large and \(p\) is small, a Poisson approximation is
also appropriate. Here,

\[
\lambda=np=3000(0.0005)=1.5.
\]

Thus,

\begin{align*}
P(X\geq2)
&\approx
1-P(Y=0)-P(Y=1)\\
&=
1-e^{-1.5}(1+1.5)\\
&\approx
0.4422,
\end{align*}

where \(Y\sim\operatorname{Pois}(1.5)\).

\end{solution}
 


%%%%%%%%%%%%%%%%%%%%%%%%%%%%%%%%%%%%%%%%%%%%%%%%%%%%%%%%%%%%%%
%%%%%%%%%%%%%%%%%%%%%%  PROBLEM 74  %%%%%%%%%%%%%%%%%%%%%%%%%%
%%%%%%%%%%%%%%%%%%%%%%%%%%%%%%%%%%%%%%%%%%%%%%%%%%%%%%%%%%%%%%

   \newpage \begin{problem}

Suppose \(X\) is a Poisson random variable satisfying

\[
P(X=1)=P(X=3).
\]

Find \(P(X=5)\).
\end{problem}
\begin{solution}

Let

\[
X\sim\operatorname{Pois}(\lambda).
\]

Then

\[
P(X=k)
=
e^{-\lambda}\frac{\lambda^k}{k!}.
\]

The condition gives

\[
e^{-\lambda}\lambda
=
e^{-\lambda}\frac{\lambda^3}{3!}.
\]

For \(\lambda>0\), canceling \(e^{-\lambda}\lambda\) yields

\[
1=\frac{\lambda^2}{6}.
\]

Therefore,

\[
\lambda^2=6
\]

and

\[
\lambda=\sqrt6.
\]

Consequently,

\begin{align*}
P(X=5)
&=
e^{-\sqrt6}
\frac{(\sqrt6)^5}{5!}\\
&=
e^{-\sqrt6}
\frac{36\sqrt6}{120}\\
&=
 {
\frac{3\sqrt6}{10}e^{-\sqrt6}.
}
\end{align*}

\end{solution}
 


%%%%%%%%%%%%%%%%%%%%%%%%%%%%%%%%%%%%%%%%%%%%%%%%%%%%%%%%%%%%%%
%%%%%%%%%%%%%%%%%%%%%%  PROBLEM 75  %%%%%%%%%%%%%%%%%%%%%%%%%%
%%%%%%%%%%%%%%%%%%%%%%%%%%%%%%%%%%%%%%%%%%%%%%%%%%%%%%%%%%%%%%
 


%%%%%%%%%%%%%%%%%%%%%%%%%%%%%%%%%%%%%%%%%%%%%%%%%%%%%%%%%%%%%%
%%%%%%%%%%%%%%%%%%%%%%  PROBLEM 76  %%%%%%%%%%%%%%%%%%%%%%%%%%
%%%%%%%%%%%%%%%%%%%%%%%%%%%%%%%%%%%%%%%%%%%%%%%%%%%%%%%%%%%%%%

   \newpage \begin{problem}

The duration \(X\) of a certain soap opera, measured in tens of hours,
has cumulative distribution function

\[
F(x)=
\begin{cases}
0, & x<4,\\[1mm]
1-\dfrac{16}{x^2}, & x\geq4.
\end{cases}
\]

\begin{enumerate}[label=(\alph*)]
    \item Find the probability density function of \(X\).
    \item Describe the graphs of \(F\) and \(f\).
    \item Find the probability that the soap opera lasts:
    \begin{enumerate}[label=(\roman*)]
        \item at most \(50\) hours;
        \item at least \(60\) hours;
        \item between \(50\) and \(70\) hours;
        \item between \(10\) and \(35\) hours.
    \end{enumerate}
\end{enumerate}
\end{problem}
\begin{solution}

\begin{enumerate}[label=(\alph*)]

\item For \(x>4\),

\[
f(x)=F'(x)
=
\frac{32}{x^3}.
\]

Thus,

\[
 {
f(x)=
\begin{cases}
\dfrac{32}{x^3}, & x\geq4,\\[2mm]
0, & x<4.
\end{cases}
}
\]

There is no point mass at \(4\), because

\[
F(4)=1-\frac{16}{16}=0.
\]

\item The CDF is zero for \(x<4\), begins at zero when \(x=4\), and
increases continuously toward \(1\).

The density is zero for \(x<4\), has value

\[
f(4)=\frac{32}{64}=\frac12,
\]

and decreases toward zero as \(x\to\infty\).

\item Since \(X\) is measured in tens of hours, \(50\) hours corresponds
to \(x=5\), \(60\) hours to \(x=6\), and so forth.

\begin{enumerate}[label=(\roman*)]

\item

\[
P(X\leq5)
=
F(5)
=
1-\frac{16}{25}
=
 {\frac9{25}}.
\]

\item Since the distribution is continuous,

\begin{align*}
P(X\geq6)
&=
1-F(6)\\
&=
\frac{16}{36}\\
&=
 {\frac49}.
\end{align*}

\item

\begin{align*}
P(5<X<7)
&=
F(7)-F(5)\\
&=
\left(1-\frac{16}{49}\right)
-
\left(1-\frac{16}{25}\right)\\
&=
\frac{16}{25}-\frac{16}{49}\\
&=
 {\frac{384}{1225}}.
\end{align*}

\item The support of \(X\) is \([4,\infty)\), corresponding to durations
of at least \(40\) hours. Hence,

\[
 {P(1<X<3.5)=0}.
\]

\end{enumerate}

\end{enumerate}

\end{solution}
 


%%%%%%%%%%%%%%%%%%%%%%%%%%%%%%%%%%%%%%%%%%%%%%%%%%%%%%%%%%%%%%
%%%%%%%%%%%%%%%%%%%%%%  PROBLEM 77  %%%%%%%%%%%%%%%%%%%%%%%%%%
%%%%%%%%%%%%%%%%%%%%%%%%%%%%%%%%%%%%%%%%%%%%%%%%%%%%%%%%%%%%%%

   \newpage \begin{problem}

The lifetime of a randomly selected used tire is \(10{,}000X\) miles,
where \(X\) has density

\[
f(x)=
\begin{cases}
\dfrac{2}{x^2}, & 1<x<2,\\[2mm]
0, & \text{otherwise}.
\end{cases}
\]

\begin{enumerate}[label=(\alph*)]
    \item What percentage of the tires last fewer than \(15{,}000\)
    miles?
    \item Among the tires lasting fewer than \(15{,}000\) miles, what
    percentage last between \(10{,}000\) and \(12{,}500\) miles?
\end{enumerate}
\end{problem}
\begin{solution}

\begin{enumerate}[label=(\alph*)]

\item A lifetime below \(15{,}000\) miles corresponds to

\[
X<1.5.
\]

Therefore,

\begin{align*}
P(X<1.5)
&=
\int_1^{3/2}\frac{2}{x^2}\,dx\\
&=
\left[-\frac2x\right]_1^{3/2}\\
&=
2-\frac43\\
&=
\frac23.
\end{align*}

Thus, $ {66\frac23\%}$ of the tires last fewer than \(15{,}000\) miles.

\item The interval \(10{,}000\) to \(12{,}500\) miles corresponds to

\[
1<X<1.25=\frac54.
\]

The required conditional probability is

\[
P\left(
1<X<\frac54
\,\middle|\,
X<\frac32
\right).
\]

Since \(\{1<X<5/4\}\subseteq\{X<3/2\}\),

\begin{align*}
P\left(
1<X<\frac54
\,\middle|\,
X<\frac32
\right)
&=
\frac{
P(1<X<5/4)
}{
P(X<3/2)
}.
\end{align*}

Now,

\begin{align*}
P\left(1<X<\frac54\right)
&=
\int_1^{5/4}\frac2{x^2}\,dx\\
&=
2-\frac{8}{5}\\
&=
\frac25.
\end{align*}

Therefore, $\frac{2/5}{2/3}
=
\frac35.$ Thus, $ {60\%}$ of the tires lasting fewer than \(15{,}000\) miles last between
\(10{,}000\) and \(12{,}500\) miles.

\end{enumerate}

\end{solution}
 


%%%%%%%%%%%%%%%%%%%%%%%%%%%%%%%%%%%%%%%%%%%%%%%%%%%%%%%%%%%%%%
%%%%%%%%%%%%%%%%%%%%%%  PROBLEM 78  %%%%%%%%%%%%%%%%%%%%%%%%%%
%%%%%%%%%%%%%%%%%%%%%%%%%%%%%%%%%%%%%%%%%%%%%%%%%%%%%%%%%%%%%%

   \newpage \begin{problem}

Let \(X\) have CDF

\[
F(x)=
\begin{cases}
0, & x<4,\\[1mm]
1-\dfrac{16}{x^2}, & x\geq4.
\end{cases}
\]

\begin{enumerate}[label=(\alph*)]
    \item Find \(E[X]\).
    \item Show that \(\operatorname{Var}(X)\) does not exist as a finite
    number.
\end{enumerate}
\end{problem}
\begin{solution}

The density is

\[
f(x)=
\frac{32}{x^3},
\qquad x\geq4.
\]

\begin{enumerate}[label=(\alph*)]

\item

\begin{align*}
E[X]
&=
\int_4^\infty x\frac{32}{x^3}\,dx\\
&=
32\int_4^\infty\frac1{x^2}\,dx\\
&=
32\left[-\frac1x\right]_4^\infty\\
&=
32\left(\frac14\right)\\
&=
 {8}.
\end{align*}

Because \(X\) is measured in tens of hours, the expected duration is

\[
 {80\text{ hours}}.
\]

\item

\begin{align*}
E[X^2]
&=
\int_4^\infty x^2\frac{32}{x^3}\,dx\\
&=
32\int_4^\infty\frac1x\,dx.
\end{align*}

The integral diverges, so

\[
E[X^2]=\infty.
\]

Therefore,

\[
 {\operatorname{Var}(X)\text{ is not finite}.}
\]

Although \(E[X]\) exists, the second moment and variance do not.

\end{enumerate}

\end{solution}
 


%%%%%%%%%%%%%%%%%%%%%%%%%%%%%%%%%%%%%%%%%%%%%%%%%%%%%%%%%%%%%%
%%%%%%%%%%%%%%%%%%%%%%  PROBLEM 79  %%%%%%%%%%%%%%%%%%%%%%%%%%
%%%%%%%%%%%%%%%%%%%%%%%%%%%%%%%%%%%%%%%%%%%%%%%%%%%%%%%%%%%%%%

   \newpage \begin{problem}

The time \(X\), in hours, required for a student to finish an aptitude
test has density

\[
f(x)=
\begin{cases}
6(x-1)(2-x), & 1<x<2,\\[1mm]
0, & \text{otherwise}.
\end{cases}
\]

Find the mean and standard deviation of \(X\).
\end{problem}
\begin{solution}

The density is symmetric about \(x=3/2\), so we expect the mean to be
\(3/2\). Directly,

\begin{align*}
E[X]
&=
\int_1^2x\,6(x-1)(2-x)\,dx\\
&=
 {\frac32}.
\end{align*}

Next,

\begin{align*}
E[X^2]
&=
\int_1^2x^2\,6(x-1)(2-x)\,dx\\
&=
\frac{23}{10}.
\end{align*}

Therefore,

\begin{align*}
\operatorname{Var}(X)
&=
E[X^2]-E[X]^2\\
&=
\frac{23}{10}
-
\left(\frac32\right)^2\\
&=
\frac{23}{10}-\frac94\\
&=
\frac{46-45}{20}\\
&=
\frac1{20}.
\end{align*}

Thus, the standard deviation is

\[
 {
\sigma_X
=
\sqrt{\frac1{20}}
=
\frac1{2\sqrt5}
\approx0.2236\text{ hours}.
}
\]

The mean is

\[
 {E[X]=1.5\text{ hours}}.
\]

\end{solution}
  

 


%%%%%%%%%%%%%%%%%%%%%%%%%%%%%%%%%%%%%%%%%%%%%%%%%%%%%%%%%%%%%%
%%%%%%%%%%%%%%%%%%%%%%  PROBLEM 80  %%%%%%%%%%%%%%%%%%%%%%%%%%
%%%%%%%%%%%%%%%%%%%%%%%%%%%%%%%%%%%%%%%%%%%%%%%%%%%%%%%%%%%%%%

   \newpage \begin{problem}

Let \(X\) be a continuous random variable with density

\[
f(x)=
\begin{cases}
\dfrac{2}{x^2}, & 1<x<2,\\[2mm]
0, & \text{otherwise}.
\end{cases}
\]

Find \(E[\ln X]\).
\end{problem}
\begin{solution}

By definition,

\[
E[\ln X]
=
\int_{1}^{2}\ln(x)\frac{2}{x^2}\,dx.
\]

Integrate by parts. Let

\[
u=\ln x,
\qquad
dv=\frac{2}{x^2}\,dx.
\]

Then

\[
du=\frac1x\,dx,
\qquad
v=-\frac2x.
\]

Therefore,

\begin{align*}
E[\ln X]
&=
\left[
-\frac{2\ln x}{x}
\right]_{1}^{2}
+
2\int_{1}^{2}\frac{1}{x^2}\,dx\\
&=
-\ln 2
+
2\left[-\frac1x\right]_{1}^{2}\\
&=
-\ln 2
+
2\left(1-\frac12\right)\\
&=
 {1-\ln 2}.
\end{align*}

Numerically,

\[
E[\ln X]\approx0.3069.
\]

\end{solution}
 


%%%%%%%%%%%%%%%%%%%%%%%%%%%%%%%%%%%%%%%%%%%%%%%%%%%%%%%%%%%%%%
%%%%%%%%%%%%%%%%%%%%%%  PROBLEM 81  %%%%%%%%%%%%%%%%%%%%%%%%%%
%%%%%%%%%%%%%%%%%%%%%%%%%%%%%%%%%%%%%%%%%%%%%%%%%%%%%%%%%%%%%%

   \newpage \begin{problem}

Let \(X\) have density

\[
f(x)=\frac12e^{-|x|},
\qquad -\infty<x<\infty.
\]

Calculate \(\operatorname{Var}(X)\).
\end{problem}
\begin{solution}

The density is symmetric about zero, so

\[
E[X]=0.
\]

Therefore,

\[
\operatorname{Var}(X)=E[X^2].
\]

Since \(x^2e^{-|x|}\) is an even function,

\begin{align*}
E[X^2]
&=
\frac12\int_{-\infty}^{\infty}x^2e^{-|x|}\,dx\\
&=
\int_{0}^{\infty}x^2e^{-x}\,dx.
\end{align*}

Using

\[
\int_{0}^{\infty}x^ne^{-x}\,dx=n!,
\]

we obtain

\[
E[X^2]=2!=2.
\]

Hence,

\[
 {\operatorname{Var}(X)=2}.
\]

\end{solution}
 


%%%%%%%%%%%%%%%%%%%%%%%%%%%%%%%%%%%%%%%%%%%%%%%%%%%%%%%%%%%%%%
%%%%%%%%%%%%%%%%%%%%%%  PROBLEM 82  %%%%%%%%%%%%%%%%%%%%%%%%%%
%%%%%%%%%%%%%%%%%%%%%%%%%%%%%%%%%%%%%%%%%%%%%%%%%%%%%%%%%%%%%%
 


%%%%%%%%%%%%%%%%%%%%%%%%%%%%%%%%%%%%%%%%%%%%%%%%%%%%%%%%%%%%%%
%%%%%%%%%%%%%%%%%%%%%%  PROBLEM 83  %%%%%%%%%%%%%%%%%%%%%%%%%%
%%%%%%%%%%%%%%%%%%%%%%%%%%%%%%%%%%%%%%%%%%%%%%%%%%%%%%%%%%%%%%

   \newpage \begin{problem}

Let \(X\) have a gamma distribution with shape parameter \(r\) and rate
parameter \(\lambda\), so that

\[
f(x)
=
\frac{\lambda^r}{\Gamma(r)}
x^{r-1}e^{-\lambda x},
\qquad x>0.
\]

Assume \(r>1\). Show that the density has a unique maximum at $\frac{r-1}{\lambda}.$
\end{problem}
\begin{solution}

Because the multiplicative constant

\[
\frac{\lambda^r}{\Gamma(r)}
\]

is positive, it is enough to maximize

\[
g(x)=x^{r-1}e^{-\lambda x}.
\]

Take logarithms:

\[
\log g(x)
=
(r-1)\log x-\lambda x.
\]

Differentiate:

\[
\frac{d}{dx}\log g(x)
=
\frac{r-1}{x}-\lambda.
\]

The critical point satisfies

\[
\frac{r-1}{x}-\lambda=0,
\]

so

\[
x=\frac{r-1}{\lambda}.
\]

Moreover,

\[
\frac{d^2}{dx^2}\log g(x)
=
-\frac{r-1}{x^2}<0
\]

for \(x>0\). Thus \(\log g\), and hence \(g\), is strictly concave at the
critical point.

Also,

\[
\frac{r-1}{x}-\lambda>0
\]

when \(x<(r-1)/\lambda\), and it is negative when
\(x>(r-1)/\lambda\). Therefore, the density increases before this point
and decreases afterward.

Hence the unique mode is

\[
 {\frac{r-1}{\lambda}}.
\]

\end{solution}
 


%%%%%%%%%%%%%%%%%%%%%%%%%%%%%%%%%%%%%%%%%%%%%%%%%%%%%%%%%%%%%%
%%%%%%%%%%%%%%%%%%%%%%  PROBLEM 84  %%%%%%%%%%%%%%%%%%%%%%%%%%
%%%%%%%%%%%%%%%%%%%%%%%%%%%%%%%%%%%%%%%%%%%%%%%%%%%%%%%%%%%%%%

   \newpage \begin{problem}
\textbf{* }Let \(X\) have a gamma distribution with shape parameter \(r\) and rate
parameter \(\lambda\). Let \(c>0\), and define

\[
Y=cX.
\]

Find the distribution function and density of \(Y\).
\end{problem}
\begin{solution}

For \(y\leq0\),

\[
F_Y(y)=0.
\]

For \(y>0\),

\begin{align*}
F_Y(y)
&=
P(Y\leq y)\\
&=
P(cX\leq y)\\
&=
P\left(X\leq\frac{y}{c}\right)\\
&=
F_X\left(\frac{y}{c}\right).
\end{align*}

Thus,

\[
 {
F_Y(y)=
\begin{cases}
0, & y\leq0,\\[1mm]
F_X\left(\dfrac{y}{c}\right), & y>0.
\end{cases}
}
\]

Using the change-of-variables formula,

\begin{align*}
f_Y(y)
&=
\frac1c
f_X\left(\frac{y}{c}\right)\\
&=
\frac1c
\frac{\lambda^r}{\Gamma(r)}
\left(\frac{y}{c}\right)^{r-1}
e^{-\lambda y/c}\\
&=
\frac{(\lambda/c)^r}{\Gamma(r)}
y^{r-1}e^{-(\lambda/c)y},
\qquad y>0.
\end{align*}

Therefore,

\[
 {
Y\sim\operatorname{Gamma}
\left(r,\frac{\lambda}{c}\right)
}
\]

under the shape--rate parameterization.

\end{solution}
 


%%%%%%%%%%%%%%%%%%%%%%%%%%%%%%%%%%%%%%%%%%%%%%%%%%%%%%%%%%%%%%
%%%%%%%%%%%%%%%%%%%%%%  PROBLEM 85  %%%%%%%%%%%%%%%%%%%%%%%%%%
%%%%%%%%%%%%%%%%%%%%%%%%%%%%%%%%%%%%%%%%%%%%%%%%%%%%%%%%%%%%%%
  


%%%%%%%%%%%%%%%%%%%%%%%%%%%%%%%%%%%%%%%%%%%%%%%%%%%%%%%%%%%%%%
%%%%%%%%%%%%%%%%%%%%%%  PROBLEM 86  %%%%%%%%%%%%%%%%%%%%%%%%%%
%%%%%%%%%%%%%%%%%%%%%%%%%%%%%%%%%%%%%%%%%%%%%%%%%%%%%%%%%%%%%%

   \newpage \begin{problem}

Let

\[
f(x)
=
\begin{cases}
\dfrac{\lambda^r}{\Gamma(r)}
x^{r-1}e^{-\lambda x},
& x>0,\\[3mm]
0,
& x\leq0,
\end{cases}
\]

where \(r>0\) and \(\lambda>0\). Prove that

\[
\int_{-\infty}^{\infty}f(x)\,dx=1.
\]
\end{problem}
\begin{solution}

Because \(f(x)=0\) for \(x\leq0\),

\[
\int_{-\infty}^{\infty}f(x)\,dx
=
\frac{\lambda^r}{\Gamma(r)}
\int_{0}^{\infty}
x^{r-1}e^{-\lambda x}\,dx.
\]

Use the substitution

\[
u=\lambda x.
\]

Then

\[
x=\frac{u}{\lambda},
\qquad
dx=\frac{du}{\lambda}.
\]

Therefore,

\begin{align*}
\int_{-\infty}^{\infty}f(x)\,dx
&=
\frac{\lambda^r}{\Gamma(r)}
\int_{0}^{\infty}
\left(\frac{u}{\lambda}\right)^{r-1}
e^{-u}\frac{du}{\lambda}\\
&=
\frac{\lambda^r}{\Gamma(r)}
\frac{1}{\lambda^r}
\int_{0}^{\infty}
u^{r-1}e^{-u}\,du\\
&=
\frac{1}{\Gamma(r)}
\Gamma(r)\\
&=
1.
\end{align*}

Thus,

\[
 {
\int_{-\infty}^{\infty}f(x)\,dx=1.
}
\]

\end{solution}
 


%%%%%%%%%%%%%%%%%%%%%%%%%%%%%%%%%%%%%%%%%%%%%%%%%%%%%%%%%%%%%%
%%%%%%%%%%%%%%%%%%%%%%  PROBLEM 87  %%%%%%%%%%%%%%%%%%%%%%%%%%
%%%%%%%%%%%%%%%%%%%%%%%%%%%%%%%%%%%%%%%%%%%%%%%%%%%%%%%%%%%%%%

   \newpage \begin{problem}

\textbf{* }Consider

\[
f(x)=
\begin{cases}
12x(1-x)^2, & 0<x<1,\\[1mm]
0, & \text{otherwise}.
\end{cases}
\]

Determine whether \(f\) is the density of a beta random variable. If so,
find \(E[X]\) and \(\operatorname{Var}(X)\).
\end{problem}
\begin{solution}

A beta density with parameters \(\alpha,\beta>0\) has the form

\[
f(x)
=
\frac{1}{B(\alpha,\beta)}
x^{\alpha-1}(1-x)^{\beta-1},
\qquad 0<x<1.
\]

Here,

\[
x(1-x)^2
=
x^{2-1}(1-x)^{3-1},
\]

so the possible parameters are

\[
\alpha=2,
\qquad
\beta=3.
\]

The beta normalizing constant is

\[
\frac{1}{B(2,3)}
=
\frac{\Gamma(5)}{\Gamma(2)\Gamma(3)}
=
\frac{4!}{1!\,2!}
=
12.
\]

Thus \(f\) is exactly the density of a

\[
\operatorname{Beta}(2,3)
\]

random variable.

For a beta random variable,

\[
E[X]=\frac{\alpha}{\alpha+\beta},
\]

so

\[
 {
E[X]=\frac25.
}
\]

Also,

\[
\operatorname{Var}(X)
=
\frac{\alpha\beta}
{(\alpha+\beta)^2(\alpha+\beta+1)}.
\]

Therefore,

\begin{align*}
\operatorname{Var}(X)
&=
\frac{(2)(3)}
{5^2(6)}\\
&=
 {\frac1{25}}.
\end{align*}

\end{solution}
 


%%%%%%%%%%%%%%%%%%%%%%%%%%%%%%%%%%%%%%%%%%%%%%%%%%%%%%%%%%%%%%
%%%%%%%%%%%%%%%%%%%%%%  PROBLEM 88  %%%%%%%%%%%%%%%%%%%%%%%%%%
%%%%%%%%%%%%%%%%%%%%%%%%%%%%%%%%%%%%%%%%%%%%%%%%%%%%%%%%%%%%%%

   \newpage \begin{problem}

\textbf{* }Determine whether

\[
f(x)=
\begin{cases}
120x^2(1-x)^4, & 0<x<1,\\[1mm]
0, & \text{otherwise}
\end{cases}
\]

is a probability density function.
\end{problem}
\begin{solution}

The function is nonnegative. It remains to check whether it integrates to
one.

The kernel

\[
x^2(1-x)^4
=
x^{3-1}(1-x)^{5-1}
\]

corresponds to a beta distribution with parameters

\[
\alpha=3,
\qquad
\beta=5.
\]

Its required normalizing constant is

\begin{align*}
\frac{1}{B(3,5)}
&=
\frac{\Gamma(8)}
{\Gamma(3)\Gamma(5)}\\
&=
\frac{7!}{2!\,4!}\\
&=
105.
\end{align*}

The proposed coefficient is \(120\), not \(105\).

Indeed,

\begin{align*}
\int_{0}^{1}120x^2(1-x)^4\,dx
&=
120B(3,5)\\
&=
\frac{120}{105}\\
&=
\frac87.
\end{align*}

Since

\[
\frac87\neq1,
\]

the proposed function is not a probability density function.

Thus,

\[
 {
f\text{ is not a probability density function.}
}
\]

Replacing \(120\) by \(105\) would produce the density of a
\(\operatorname{Beta}(3,5)\) random variable.

\end{solution}






%%%%%%%%%%%%%%%%%%%%%%%%%%%%%%%%%%%%%%%%%%%%%%%%%%%%%%%%%%%%%%%%%%%%%%%%%%%%%%%
% BINOMIAL AND DISTRIBUTION-APPROXIMATION PROBLEMS
%%%%%%%%%%%%%%%%%%%%%%%%%%%%%%%%%%%%%%%%%%%%%%%%%%%%%%%%%%%%%%%%%%%%%%%%%%%%%%%


\newpage
\begin{problem}

Suppose that
\[
X\sim\operatorname{Bin}(100,0.4).
\]

Use a Normal approximation, with an appropriate continuity correction, to
approximate
\[
\mathbb{P}(35\leq X\leq45).
\]

Clearly verify that the conditions for the Normal approximation are satisfied
and state the approximating Normal distribution.

\end{problem}

\begin{solution}

We have
\[
X\sim\operatorname{Bin}(100,0.4).
\]

To check whether the Normal approximation is appropriate, we compute
\[
np=100(0.4)=40
\]
and
\[
n(1-p)=100(0.6)=60.
\]

Since both quantities are greater than \(10\), the Normal approximation is
appropriate.

The mean and variance are
\[
\mu=np=40
\]
and
\[
\sigma^2=np(1-p)=100(0.4)(0.6)=24.
\]

Therefore, we approximate \(X\) by
\[
Y\sim N(40,24).
\]

Using the continuity correction,
\[
\mathbb{P}(35\leq X\leq45)
\approx
\mathbb{P}(34.5<Y<45.5).
\]

Standardizing,
\[
\begin{aligned}
\mathbb{P}(34.5<Y<45.5)
&=
\mathbb{P}\left(
\frac{34.5-40}{\sqrt{24}}
<
Z
<
\frac{45.5-40}{\sqrt{24}}
\right)\\
&=
\mathbb{P}(-1.123<Z<1.123).
\end{aligned}
\]

Therefore,
\[
\begin{aligned}
\mathbb{P}(35\leq X\leq45)
&\approx
\Phi(1.123)-\Phi(-1.123)\\
&=
2\Phi(1.123)-1\\
&\approx
 {0.739}.
\end{aligned}
\]

\end{solution}



%%%%%%%%%%%%%%%%%%%%%%%%%%%%%%%%%%%%%%%%%%%%%%%%%%%%%%%%%%%%%%%%%%%%%%%%%%%%%%%


\newpage
\begin{problem}

The odds are \(1\) to \(5000\) in favor of a customer buying a particular
fiction bestseller.

Suppose that \(800\) customers enter a bookstore each day and that a month
has \(30\) days.

How many copies of the book should the bookstore stock each month so that,
with probability greater than \(98\%\), it does not run out?

Determine whether a Poisson approximation or a Normal approximation is more
appropriate, and verify the required conditions before solving the problem.

\end{problem}

\begin{solution}

Let
\[
X=\text{the number of customers who buy the book during the month}.
\]

The total number of customers during the month is
\[
n=30(800)=24000.
\]

Because the odds in favor of buying the book are \(1\) to \(5000\), the
probability that a customer buys the book is
\[
p=\frac{1}{1+5000}=\frac{1}{5001}.
\]

Therefore,
\[
X\sim\operatorname{Bin}\left(24000,\frac{1}{5001}\right).
\]

We first determine the appropriate approximation.

The expected number of successes is
\[
np
=
24000\left(\frac{1}{5001}\right)
\approx4.799.
\]

Although \(n\) is large,
\[
np\approx4.799<10,
\]
so the usual Normal approximation is not appropriate.

However, \(n=24000\) is large and
\[
p=\frac{1}{5001}
\]
is very small. Therefore, the Poisson approximation is appropriate, with
\[
\lambda=np=\frac{24000}{5001}\approx4.799.
\]

Thus,
\[
X\approx W,
\qquad
W\sim\operatorname{Pois}(4.799).
\]

If the bookstore stocks \(m\) copies, then it does not run out when
\[
X\leq m.
\]

We seek the smallest integer \(m\) such that
\[
\mathbb{P}(X\leq m)>0.98.
\]

Using the Poisson approximation,
\[
\mathbb{P}(X\leq9)
\approx
\mathbb{P}(W\leq9)
=
\sum_{k=0}^{9}
e^{-4.799}\frac{4.799^k}{k!}
\approx0.9749.
\]

Since
\[
0.9749<0.98,
\]
stocking \(9\) copies is not sufficient.

Next,
\[
\mathbb{P}(X\leq10)
\approx
\mathbb{P}(W\leq10)
=
\sum_{k=0}^{10}
e^{-4.799}\frac{4.799^k}{k!}
\approx0.9896.
\]

Since
\[
0.9896>0.98,
\]
the smallest satisfactory value is \(m=10\).

Therefore, the bookstore should stock
\[
 {10\text{ copies}.}
\]

\end{solution}



%%%%%%%%%%%%%%%%%%%%%%%%%%%%%%%%%%%%%%%%%%%%%%%%%%%%%%%%%%%%%%%%%%%%%%%%%%%%%%%


\newpage
\begin{problem}

Every day, a factory produces \(5000\) light bulbs, of which \(2500\) are
Type I and \(2500\) are Type II.

A sample of \(40\) light bulbs is selected uniformly at random without
replacement.

Approximate the probability that the sample contains at least \(18\) light
bulbs of each type.

Clearly explain and justify every approximation that you use.

\end{problem}

\begin{solution}

Let
\[
X=\text{the number of Type I bulbs in the sample}.
\]

Because the sample is selected without replacement, the exact distribution is
\[
X\sim\operatorname{Hypergeo}(5000,2500,40).
\]

The sampling fraction is
\[
\frac{40}{5000}=0.008.
\]

Since the sample size is very small relative to the population size, sampling
without replacement is approximately equivalent to sampling with replacement.
Therefore,
\[
X
\approx
\operatorname{Bin}\left(40,\frac{2500}{5000}\right)
=
\operatorname{Bin}\left(40,\frac12\right).
\]

For this Binomial distribution,
\[
np=40\left(\frac12\right)=20
\]
and
\[
n(1-p)=40\left(\frac12\right)=20.
\]

Since both quantities are greater than \(10\), the Normal approximation is
appropriate.

The mean and variance are
\[
\mu=np=20
\]
and
\[
\sigma^2=np(1-p)
=
40\left(\frac12\right)\left(\frac12\right)
=
10.
\]

Thus, we further approximate \(X\) by
\[
Y\sim N(20,10).
\]

If there are \(X\) Type I bulbs, then there are
\[
40-X
\]
Type II bulbs.

The sample contains at least \(18\) bulbs of each type when
\[
X\geq18
\]
and
\[
40-X\geq18.
\]

The second inequality is equivalent to
\[
X\leq22.
\]

Therefore, the desired event is
\[
18\leq X\leq22.
\]

Using the continuity correction,
\[
\mathbb{P}(18\leq X\leq22)
\approx
\mathbb{P}(17.5<Y<22.5).
\]

Standardizing,
\[
\begin{aligned}
\mathbb{P}(17.5<Y<22.5)
&=
\mathbb{P}\left(
\frac{17.5-20}{\sqrt{10}}
<
Z
<
\frac{22.5-20}{\sqrt{10}}
\right)\\
&=
\mathbb{P}(-0.791<Z<0.791).
\end{aligned}
\]

Therefore,
\[
\begin{aligned}
\mathbb{P}(18\leq X\leq22)
&\approx
\Phi(0.791)-\Phi(-0.791)\\
&=
2\Phi(0.791)-1\\
&\approx
2(0.785)-1\\
&=
 {0.570}.
\end{aligned}
\]

Thus, the approximate probability that the sample contains at least \(18\)
bulbs of each type is \(0.570\).

\end{solution}



%%%%%%%%%%%%%%%%%%%%%%%%%%%%%%%%%%%%%%%%%%%%%%%%%%%%%%%%%%%%%%%%%%%%%%%%%%%%%%%




%%%%%%%%%%%%%%%%%%%%%%%%%%%%%%%%%%%%%%%%%%%%%%%%%%%%%%%%%%%%%%%%%%%%%%%%%%%%%%%
% DIRECT NORMAL-DISTRIBUTION PROBLEMS
%%%%%%%%%%%%%%%%%%%%%%%%%%%%%%%%%%%%%%%%%%%%%%%%%%%%%%%%%%%%%%%%%%%%%%%%%%%%%%%


\newpage
\begin{problem}

Suppose that the lifetimes of light bulbs produced by a company are Normally
distributed with mean \(1000\) hours and standard deviation \(100\) hours.

The company claims that \(95\%\) of its light bulbs last at least \(900\)
hours.

Is the company's claim correct? Justify your answer mathematically.

\end{problem}

\begin{solution}

Let
\[
X=\text{the lifetime of a randomly selected light bulb}.
\]

Then
\[
X\sim N(1000,100^2).
\]

The proportion of bulbs that last at least \(900\) hours is
\[
\mathbb{P}(X\geq900).
\]

Standardizing,
\[
\begin{aligned}
\mathbb{P}(X\geq900)
&=
\mathbb{P}\left(
Z\geq\frac{900-1000}{100}
\right)\\
&=
\mathbb{P}(Z\geq-1).
\end{aligned}
\]

By symmetry of the standard Normal distribution,
\[
\mathbb{P}(Z\geq-1)
=
\mathbb{P}(Z\leq1)
=
\Phi(1).
\]

Using the standard Normal table,
\[
\Phi(1)\approx0.8413.
\]

Therefore,
\[
\mathbb{P}(X\geq900)\approx0.8413.
\]

Thus, approximately \(84.13\%\), rather than \(95\%\), of the bulbs last at
least \(900\) hours. Since
\[
0.8413<0.95,
\]
the company's claim is incorrect.

\[
 {\text{No, the company's claim is not correct.}}
\]

\end{solution}



%%%%%%%%%%%%%%%%%%%%%%%%%%%%%%%%%%%%%%%%%%%%%%%%%%%%%%%%%%%%%%%%%%%%%%%%%%%%%%%


\newpage
\begin{problem}

The lifetime of a bulb produced by the first company is Normally distributed
with mean \(1000\) hours and standard deviation \(100\) hours.

The lifetime of a bulb produced by the second company is Normally distributed
with mean \(900\) hours and standard deviation \(150\) hours.

Howard buys one bulb from each company. Assuming that the two bulb lifetimes
are independent, what is the probability that at least one of the bulbs lasts
\(980\) hours or more?

\end{problem}

\begin{solution}

Let
\[
X=\text{the lifetime of the bulb from the first company}
\]
and
\[
Y=\text{the lifetime of the bulb from the second company}.
\]

Then
\[
X\sim N(1000,100^2)
\]
and
\[
Y\sim N(900,150^2).
\]

We want
\[
\mathbb{P}(X\geq980\text{ or }Y\geq980).
\]

It is easier to use the complementary event:
\[
\mathbb{P}(X\geq980\text{ or }Y\geq980)
=
1-\mathbb{P}(X<980,\ Y<980).
\]

Because \(X\) and \(Y\) are independent,
\[
\mathbb{P}(X<980,\ Y<980)
=
\mathbb{P}(X<980)\mathbb{P}(Y<980).
\]

For the first bulb,
\[
\begin{aligned}
\mathbb{P}(X<980)
&=
\mathbb{P}\left(
Z<
\frac{980-1000}{100}
\right)\\
&=
\mathbb{P}(Z<-0.20)\\
&=
\Phi(-0.20)\\
&\approx0.4207.
\end{aligned}
\]

For the second bulb,
\[
\begin{aligned}
\mathbb{P}(Y<980)
&=
\mathbb{P}\left(
Z<
\frac{980-900}{150}
\right)\\
&=
\mathbb{P}(Z<0.533)\\
&=
\Phi(0.533)\\
&\approx0.7031.
\end{aligned}
\]

Therefore,
\[
\begin{aligned}
\mathbb{P}(X<980,\ Y<980)
&\approx
(0.4207)(0.7031)\\
&\approx
0.2958.
\end{aligned}
\]

Consequently,
\[
\begin{aligned}
\mathbb{P}(X\geq980\text{ or }Y\geq980)
&\approx
1-0.2958\\
&=
 {0.7042}.
\end{aligned}
\]

Thus, the probability that at least one of the two bulbs lasts \(980\) hours
or more is approximately \(70.42\%\).

\end{solution}


















 

\newpage
\section*{References}

\begin{enumerate}
    \item Sheldon Ross,
    \textit{A First Course in Probability},
    10th ed.,
    Pearson, 2019.

    \item Saeed Ghahramani,
    \textit{Fundamentals of Probability with Stochastic Processes},
    4th ed.,
    Pearson, 2018.


        \item Arman Jahangiri,
    \textit{MATH 394: Probability I --- Lecture Notes},
    Department of Mathematics,
    University of Washington,
    Summer 2026.

    \item Arman Jahangiri,
    \textit{MATH 394: Probability I --- Homework Problems and Solutions},
    Department of Mathematics,
    University of Washington,
    Summer 2026.
\end{enumerate}















\end{document}
